\chapter{Homomorfismos y construccion de cuerpos}

Nuestro proximo objetivo es buscar cuerpos donde un polinomio dado tenga todas sus raíces.

\section{Homomorfismos de cuerpos. Cuerpos de descomposicion}



\begin{lema}
    Sea $F$ un cuerpo y $I$ un ideal suyo, entonces $I = \{0\}$ o $I = F$.
    \begin{proof}
        Supuesto que $I\neq \{0\}$, existe por tanto $a\in I\setminus\{0\}$. Sea $b\in F$, tenemos que:
        \begin{equation*}
            b = b\cdot a^{-1}\cdot a
        \end{equation*}
        Por lo que $b\in I$, de donde $I = F$.
    \end{proof}
\end{lema}

\begin{lema}
    Sea $\sigma:F\to A$ un homomorfismo de anillos donde $F$ es un cuerpo y $A$ es no trivial. Entonces $\sigma$ es inyectivo y, por tanto, $Im \sigma$ es un cuerpo isomorfo a $F$ y subanillo de $A$.
    \begin{proof}
        Solo hemos de probar que $\ker \sigma = \{0\}$. Para ello, $\ker \sigma$ es un ideal de $F$ que no es $F$ (ya que $\sigma(1) = 1\neq 0$), de donde $\ker \sigma = \{0\}$. Para ver que $Im \sigma\cong F$, basta aplicar el Primer Teorema de Isomorfía:
        \begin{equation*}
            F = \dfrac{F}{\ker \sigma} \cong Im \sigma
        \end{equation*}
        Por ser $\sigma$ un homomorfismo de anillos tenemos que $Im\sigma$ es subanillo de $A$.
    \end{proof}
\end{lema}

\begin{definicion}[Homomorfismo de cuerpos]
    Sea $F\stackrel{\sigma}{\to} K$ un homomorfismo de anillos entre cuerpos, diremos entonces que es un \underline{homomorfismo de cuerpos}.
\end{definicion}

\begin{observacion}
    Resulta sorprendente que exigir ``buenas propiedades'' a una aplicación entre anillos ya nos da una aplicación con ``buenas propiedades'' entre cuerpos, pero resulta que lo único que nos faltaba era que la aplicación se comporte bien con los inversos, propiedad que queda garantizada al exigir ``buenas propiedades'' sobre anillos:
    \begin{equation*}
        1 = \sigma(1) = \sigma\left(\alpha\alpha^{-1}\right) = \sigma(\alpha)\sigma\left(\alpha^{-1}\right) \Longrightarrow \sigma\left(\alpha^{-1}\right) = {\sigma(\alpha)}^{-1}
    \end{equation*}
    Como por el Lema anterior todo homomorfismo de cuerpos $F\stackrel{\sigma}{\to}K$ es siempre inyectivo, tendremos siempre una copia de $F$ dentro de $K$, que en ocasiones identificaremos con el propio $F$, viendo $\sigma(F)$ como una copia isomorfa de $F$. Como $\sigma(F)\leq K$ es una extensión de cuerpos, podemos ver $K$ como un $\sigma(F)-$espacio vectorial. Además, si identificamos $F$ con $\sigma(F)$, podremos ver $K$ como un $F-$espacio vectorial.
\end{observacion}

\begin{definicion}
    Siempre que tengamos $F\stackrel{\sigma}{\to} K$ y $f\in F[x]$ dada por:
\begin{equation*}
    f = \sum_{i=0}^{n} f_i x^i, \qquad f_i \in F \quad \forall i \in \{1,\ldots,n\}
\end{equation*}
Definiremos:
\begin{equation*}
    f^\sigma = \sum_{i=0}^{n} \sigma(f_i) x^i \in K[x]
\end{equation*}
Se verifica que la correspondencia $f\mapsto f^\sigma$ es un homomorfismo de anillos entre $F[x]$ y $K[x]$, por la Propiedad Universal del anillo de polinomios.
\end{definicion}

\begin{ejemplo}\label{ejemplo:cuerpo_descomposicion}
    Sea $f\in F[x]$, $f$ no constante, sea $p\in F[x]$ un factor irreducible de $f$, consideramos\footnote{Donde $\langle p \rangle $ es el ideal generado por $p$.}:
    \begin{equation*}
        K = \dfrac{F[x]}{\langle p \rangle }
    \end{equation*}
    
    Como $p$ es irreducible, tenemos que $K$ es un cuerpo. Definimos $\sigma:F\to K$ como:
    \begin{equation*}
        \sigma(a) = a + \langle p \rangle  \qquad \forall a\in F
    \end{equation*}
    que es un homomorfismo de anillos como composición de la inclusión en $F[x]$ con la proyección al cociente:
    \begin{figure}[H]
        \centering
        \shorthandoff{""}
        \begin{tikzcd}
            F \arrow[r, "\iota", hook] & {F[x]} \arrow[r, "\pi"] & {\dfrac{F[x]}{\langle p\rangle}}
        \end{tikzcd}
        \shorthandon{""}
    \end{figure}
    \noindent
    Por tanto, se trata de un homomorfismo de cuerpos, con:
    \begin{equation*}
        \sigma(F) = \{a+\langle p \rangle : a\in F\} \cong F
    \end{equation*}
    Sea $\alpha = x+\langle p \rangle \in K$, tenemos que:
    \begin{equation*}
        p^\sigma(\alpha) = \sum_{i=0}^{n} (p_i + \langle p \rangle){(x+\langle p \rangle )}^{i} = \sum_{i=0}^{n} p_ix^i + \langle p \rangle  = p + \langle p \rangle  \AstIg 0 + \langle p \rangle 
    \end{equation*}
    donde en $(\ast)$ he empleado que $p|p$.
    Además:
    \begin{equation*}
        \sigma(F)(\alpha) = K
    \end{equation*}
    \begin{description}
        \item [$\subseteq )$] Basta ver que $K$ contiene a $\sigma(F)$ y a $\alpha$.
        \item [$\supseteq )$] Si tomamos un elemento de $K$, este será de la forma $g+\langle p \rangle $ para cierta $g\in F[x]$ dada por:
            \begin{equation*}
                \sum_{i=0}^{n}g_ix^i, \qquad g_i \in F, \quad \forall i \in \{1,\ldots,n\}
            \end{equation*}
            Por lo que:
            \begin{equation*}
                g+\langle p \rangle  = \sum_{i=0}^{n}g_ix^i + \langle p \rangle  = \sum_{i=0}^{n}(g_i+\langle p \rangle ){(x+\langle p \rangle )}^{i}
                = \sum_{i=0}^{n}(g_i+\langle p \rangle ){\alpha}^{i} \in \sigma(F)(\alpha)
            \end{equation*}
    \end{description}
\end{ejemplo}


El ejemplo anterior es de gran importancia, pues nos explica que dado cualquier cuerpo $F$ y una ecuación polinómica que no sabemos resolver en $F$ (o equivalentemente, un polinomio irreducible $p$ de grado mayor\footnote{Notemos que en un cuerpo siempre sabemos resolver una ecuación del tipo $ax=b$, las correspondientes a los polinomios de grado 1.} o igual que dos), podemos encontrar siempre un cuerpo más grande $K$ (que contiene una copia isomorfa de $F$) en el que hayamos una solución de dicha ecuación (como el elemento $x+\langle p \rangle $).

\begin{comment}
\begin{ejemplo}
\noindent Consideremos el cuerpo $F = \mathbb{R}$ y el polinomio $p(x) = x^2 + 1 \in \mathbb{R}[x]$, el cual es irreducible sobre $\mathbb{R}$ al no poseer raíces reales. Definimos el cuerpo cociente:
\[
K = \frac{\mathbb{R}[x]}{\langle x^2 + 1 \rangle}
\]

El lector podrá pensar que $K$ debería ser $\bb{C}$.
    Definimos el homomorfismo de evaluación en la unidad imaginaria $i \in \mathbb{C}$:
    \Func{\psi}{\mathbb{R}[x]}{\mathbb{C}}{g(x)}{g(i)}
    Dicha aplicación satisface:
    \begin{itemize}
        \item Es sobreyectiva, pues dado cualquier $a + bi \in \mathbb{C}$, se tiene $\psi(a + bx) = a + bi$.
        \item Su núcleo es:
        \[
        \ker(\psi) = \{g(x) \in \mathbb{R}[x] \mid g(i) = 0\} = \langle x^2 + 1 \rangle
        \]
        dado que $x^2 + 1$ es el polinomio mónico de grado mínimo en $\mathbb{R}[x]$ que anula a $i$.
    \end{itemize}
    Aplicando el Primer Teorema de Isomorfía para anillos, inducimos un isomorfismo de cuerpos:
    \[
    \bar{\psi} \colon \frac{\mathbb{R}[x]}{\langle x^2 + 1 \rangle} \overset{\cong}{\longrightarrow} \mathbb{C}
    \]
    que identifica la clase $x + \langle x^2 + 1 \rangle$ con la unidad imaginaria $i$. Concluimos así que $K \cong \mathbb{C}$.
\end{ejemplo}
\end{comment}

\begin{notacion}
    Siempre que estemos trabajando con un cuerpo $F$ y digamos que ``existe un homomorfismo $F\stackrel{\sigma}{\to}K$'', lo que queremos decir en realidad es que existen otro cuerpo $K$ y un homomorfismo de cuerpos entre ellos $\sigma:F\to K$, pero usaremos la primera expresión para abreviar.
\end{notacion}

\begin{lema}\label{lema:cuerpo_descomposicion}
    Si $f\in F[x]$ es no constante y $p$ es un factor irreducible de $f$, entonces existen $F\stackrel{\sigma}{\to}K$ homomorfismo de cuerpos y $\alpha\in K$ tales que:
    % \begin{equation*}
    %     f^\sigma(\alpha) = 0 \qquad \text{y}\qquad  \sigma(F)(\alpha) = K
    % \end{equation*}
    \begin{equation*}
        p^\sigma(\alpha) = 0 \qquad \text{y}\qquad  K = \sigma(F)(\alpha)
        \qquad \text{y}\qquad [K:\sigma(F)] = \deg p
    \end{equation*}
    Bajo estas condiciones, a menudo identificaremos $F$ con $\sigma(F)$ y en dicho caso, escribiremos $K= F(\alpha)$.
    \begin{proof}
        La demostración se deduce del ejemplo de la página~\pageref{ejemplo:cuerpo_descomposicion}. Respecto al grado de la extensión, tenemos que:
        \begin{equation*}
            [K:\sigma(F)] = [\sigma(F)(\alpha):\sigma(F)] = \deg \Irr(\alpha,\sigma(F))
        \end{equation*}

        Como $\sigma$ es un homomorfismo de cuerpos (luego inyectivo) entre $F$ y $K$, al restringir el codominio de $\sigma$ a $\sigma(F)$, tenemos que $\sigma:F\to \sigma(F)$ es un isomorfismo de cuerpos, por lo que lleva irreducibles a irreducibles, de donde $p^\sigma$ es irreducible en $\sigma(F)[x]$. Además, $\alpha$ es raíz de $p^\sigma$, por lo que $\Irr(\alpha,\sigma(F)) = p^\sigma$, de donde:
        \begin{equation*}
            [K:\sigma(F)] = \deg \Irr(\alpha,\sigma(F)) = \deg p^\sigma = \deg p
        \end{equation*}    
    \end{proof}
\end{lema}

\begin{prop}[Existencia de cuerpos de descomposición]\label{prop:existe_cdd}~\\
    Sea $f\in F[x]$ no constante. Entonces existe un homomorfismo de cuerpos $\sigma:F\to E$ tal que $E$ es un cuerpo de descomposición de $f^{\sigma}$.
    \begin{proof}
        Suponemos sin pérdida de generalidad que $f$ es mónico. Vamos a ver que existe $F\stackrel{\sigma}{\to}L$ tal que $f^\sigma$ se descompone completamente como producto de factores lineales en $L[x]$. Para ello, descomponemos $f=gh$, donde $g\in F[x]$ es producto de polinomios lineales y $h\in F[x]$ es un polinomio sin raíces en $F$. Por inducción sobre el grado de $h$ (usando el segundo principio de inducción):
        \begin{itemize}
            \item \underline{Si $\deg h = 0$}, tomando $L = F$ y $\sigma = \id_F$ se tiene.
            \item \underline{Supuesto que $\deg h>0$ y la hipótesis de inducción}, tomamos $p$ un factor irreducible de $h$, por lo que podemos aplicar el Lema~\ref{lema:cuerpo_descomposicion}, con lo que existen $F\stackrel{\tau}{\to}K$ y $\alpha\in K$ tal que $p^\tau(\alpha)=0$ y $K = \tau(F)(\alpha)$. Observamos que $h^\tau(\alpha)=0$.

                El polinomio $g$ que habíamos escogido será de la forma:
                \begin{equation*}
                    g = (x-\alpha_1)\cdot  \ldots \cdot (x-\alpha_t), \qquad \alpha_1,\ldots,\alpha_t \in F
                \end{equation*}
                Extraemos ahora los factores lineales de $h^\tau$ en $K[x]$ (sabemos que al menos $s\geq 1$, puesto que $\alpha$ es raíz de $h^\tau$):
                \begin{equation*}
                    h^\tau = (x-\beta_1)\cdot \ldots\cdot (x-\beta_s)k, \qquad k\in K[x], \beta_1,\ldots,\beta_s\in K
                \end{equation*}
                Con uno de los $\beta_i$ es $\alpha$ y $k$ sin raíces en $K$ y de grado menor que el de $h^\tau$. En definitiva, tenemos que:
                \begin{equation*}
                    f^\tau = g^\tau h^\tau = (x-\tau(\alpha_1))\cdot  \ldots \cdot (x-\tau(\alpha_s))(x-\beta_1) \cdot \ldots \cdot (x-\beta_s)k, \qquad deg k < deg h^\tau
                \end{equation*}
                Aplicando la hipótesis de inducción tomando $k$ como $h$, existe un homomorfismo $K\stackrel{\rho}{\to}L$ tal que $k^\rho$  se descompone como producto de polinomios lineales en $L[x]$. En definitiva, tendremos que ${(f^\tau)}^{\rho}$ se descompone como producto de polinomios lineales en $L[x]$. Si tomamos:
                \begin{equation*}
                    \sigma = \rho\circ \tau : F\to L
                \end{equation*}
                tenemos que $f^\sigma$ se descompone como producto de lineales en $L[x]$.
        \end{itemize}
        Una vez tenemos que existe $\sigma:F\to L$ de forma que $f^\sigma$ se descompone como producto de polinomios lineales en $L[x]$, si $\gamma_1,\ldots,\gamma_r\in L$ son las raíces de $f^\sigma$, podemos considerar $E=\sigma(F)(\gamma_1,\ldots,\gamma_r)\leq L$, con lo que la restricción en codominio de $\sigma$ a $E$ nos da un homomorfismo, donde $E$ es un cuerpo de descomposición de $f^\sigma$.
    \end{proof}
\end{prop}

\begin{definicion}
    Sea $f\in F[x]$, diremos que un homomorfismo de cuerpos ${\sigma:F\to E}$ es un cuerpo de descomposición de $f$ si $E$ es un cuerpo de descomposición de $f^\sigma$, o equivalentemente, si:
    \begin{itemize}
        \item $f^\sigma$ se descompone en $E[x]$ como producto de polinomios de grado 1, siendo $\alpha_1, \ldots, \alpha_t\in E$ todas sus raíces.
        \item $E = \sigma(F)(\alpha_1,\ldots,\alpha_t)$.
    \end{itemize}
\end{definicion}

\noindent
Respecto a esta última definición, debemos tener claro que en la Definición~\ref{def:cuerpo_descomposicion} llamamos cuerpo de descomposición de $f\in F[x]$ a una extensión $F\leq K$ de forma que $K$ cumplía ciertas propiedades relativas a $f$. Ahora, lo que hacemos es ver el homomorfismo $F\stackrel{\sigma}{\to} K$ como una extensión de cuerpos, identificando $F$ con $\sigma(F)$, por lo que al propio homomorfismo (que determina de qué forma $K$ es extensión de $F$) le llamamos ahora cuerpo de descomposición, si $K$ cumple unas propiedades relativas a $f^\sigma$.

\noindent
Sin embargo, no hemos cambiado el concepto, pues si tenemos $F\leq K$ con $K$ cuerpo de descomposición de $f\in F[x]$ tenemos que $\iota:F\hookrightarrow K$ es un cuerpo de descomposición de $f$. Hemos generalizado el concepto de cuerpo de descomposición.\\

\noindent
A partir de la Proposición~\ref{prop:existe_cdd}, para cualquier $f\in F[x]$ tenemos garantizada siempre la existencia de un cuerpo de descomposición de $f$.

\begin{ejemplo}
    Tomamos $f = x^2+x+1\in \bb{Z}_2[x]$. Como $f(0) = f(1) = 1 \neq 0$, tenemos que $f$ no tiene raíces en $\bb{Z}_2$. Nuestro objetivo es buscar un cuerpo de descomposición suyo.\\

    \noindent
    Observemos que, como $f$ es de grado 2 y no tiene raíces en $\bb{Z}_2$, $f$ es irreducible, por lo que 
    por el Lema~\ref{lema:cuerpo_descomposicion}, existen $\bb{Z}_2\stackrel{\sigma}{\to}K$ y $\alpha\in K$ tal que $f^\sigma(\alpha) = 0$, $K = \sigma(F)(\alpha)$ y $[K:\sigma(F)] = \deg f = 2$.\\

    Como $\sigma(\bb{Z}_2)\cong \bb{Z}_2$, tenemos que $\sigma(\bb{Z}_2)=\{0,1\}\leq K$. Identificando $\sigma(\bb{Z}_2)$ con $\bb{Z}_2$, comprobamos que $f^\sigma = x^2+x+1\in \sigma(\bb{Z}_2)[x]$ es irreducible, por lo que
    $$\Irr(\alpha,\sigma(\bb{Z}_2)) = f^\sigma
    $$

    De ahí, deducimos que $[K:\sigma(\bb{Z}_2)] = \deg \Irr(\alpha,\sigma(\bb{Z}_2)) = 2$ y que $\{1,\alpha\}$ es una $\sigma(\bb{Z}_2)-$base de $K$. Por tanto, cada elemento de $K$ es de la forma:
    \begin{equation*}
        a + b\alpha, \qquad a,b\in \sigma(\bb{Z}_2)=\{0,1\}
    \end{equation*}

    Evaluando los cuatro posibles pares de valores de $a$ y $b$, obtenemos que los elementos de $K$ son:
    \begin{equation*}
        K=\{0,1,\alpha,1+\alpha\}
    \end{equation*}

    Como $\alpha$ es raíz de $f^\sigma$, tenemos que $\alpha^2+\alpha+1=0$, por lo que haciendo uso de $\bb{Z}_2\cong \sigma(\bb{Z}_2)$ tenemos que $1+\alpha = \alpha^2$. Por tanto, podemos describir también a $K$ como:
    \begin{equation*}
        K=\{0,1,\alpha,\alpha^2\}
    \end{equation*}

    Por último, dado que $\alpha\in K$ es raíz de $f^\sigma$, podemos factorizar $f^\sigma$ en $K[x]$. Dividimos haciendo uso de $\bb{Z}_2\cong \sigma(\bb{Z}_2)$:
    \[
    \renewcommand{\arraystretch}{1.2}
    \setlength{\arraycolsep}{2pt}
    \begin{array}{rrrrr@{\quad}l}
    x^2 & + & x          & + & 1        & \multicolumn{1}{|l}{x + \alpha} \\ \cline{6-6}
    -x^2 & - & \alpha x   &   &          & \multicolumn{1}{|l}{x + \alpha^2} \\ \cline{1-3}
        &   & \alpha^2 x & + & 1        & \\
        & - & \alpha^2 x & - & \alpha^3 & \\ \cline{2-5}
        &   &            &   & 0        & 
    \end{array}
    \]
    De donde deducimos que $f^\sigma = (x+\alpha)(x+\alpha^2)$. De aquí, deducimos que un cuerpo de descomposición de $f$ es:
    $$\sigma(\bb{Z}_2)(\alpha, \alpha^2)=\sigma(\bb{Z}_2)(\alpha) = K$$
\end{ejemplo}

\begin{observacion}
    Si tenemos $F\leq K$ una extensión de cuerpos con $[K:F] = n$ y $|F| = m$ tendremos entonces que $|K| = m^n$, ya que cada elemento de $K$ queda unívocamente determinado por la elección de $n$ elementos de $F$, entre los cuales tenemos $m$ para elegir.
\end{observacion}

\begin{ejemplo}
    Al igual que en el ejemplo anterior, buscamos un cuerpo de descomopsición de:
    \begin{equation*}
        f = x^3+x+1 \in \bb{Z}_2[x]
    \end{equation*}

    Dado que $f(0) = 1 \neq 0$ y $f(1) = 1 \neq 0$, $f$ no tiene raíces en $\mathbb{Z}_2$. Al tener grado 3, deducimos que $f$ es irreducible en $\mathbb{Z}_2[x]$. Al igual que en el ejemplo anterior, aplicando el Lema~\ref{lema:cuerpo_descomposicion} obtenemos que existen $\bb{Z}_2\stackrel{\sigma}{\to}K$ y $\alpha\in K$ tal que $f^\sigma(\alpha) = 0$, $K = \sigma(F)(\alpha)$ y $[K:\sigma(F)] = \deg f = 3$.\\

    Identificando $\sigma(\mathbb{Z}_2) \cong \mathbb{Z}_2$, se tiene que $[L : \sigma(\mathbb{Z}_2)] = \deg \operatorname{Irr}(\alpha, \sigma(\mathbb{Z}_2)) = 3$. Por tanto, una $\sigma(\mathbb{Z}_2)-$base de $K$ es $\{1,\alpha,\alpha^2\}$, y cada elemento de $K$ es de la forma:
    \begin{equation*}
        a + b\alpha + c\alpha^2, \qquad a,b,c\in \sigma(\mathbb{Z}_2) = \{0,1\}
    \end{equation*}
    Por lo que $|K| = 2^3 = 8$. Puesto que $\alpha^3 + \alpha + 1 = 0$, tenemos las siguientes relaciones:
    \begin{equation*}
        \begin{array}{ccc|l}
            a & b & c & a + b\alpha + c\alpha^2 \\
            \hline
            0 & 0 & 0 & 0 \\
            1 & 0 & 0 & 1 \\
            0 & 1 & 0 & \alpha \\
            0 & 0 & 1 & \alpha^2 \\
            1 & 1 & 0 & 1 + \alpha = \alpha^3 \\
            1 & 0 & 1 & 1 + \alpha^2 =(1+\alpha)^2 = \alpha^6 \\
            0 & 1 & 1 & \alpha + \alpha^2 = \alpha(1+\alpha) = \alpha^4 \\
            1 & 1 & 1 & 1 + \alpha + \alpha^2 = \alpha^3 + \alpha^2 = \alpha^2(1+\alpha) = \alpha^5
        \end{array}
    \end{equation*}
    Por tanto, podemos describir a $K$ como:
    \begin{equation*}
        K = \{0,1,\alpha,\alpha^2,\alpha^3,\alpha^4,\alpha^5,\alpha^6\}
    \end{equation*}
    donde $\alpha$ es una raíz de $f^\sigma$.\\

    Buscamos ahora la factorización de $f^\sigma$ en $K[x]$. Como $\alpha$ es raíz de $f^\sigma$, realizamos la división de polinomios:
    \[
    \renewcommand{\arraystretch}{1.2}
    \setlength{\arraycolsep}{2pt}
    \begin{array}{rrrrrrr@{\quad}l}
    x^3 &   &      & + & x          & + & 1                  & \multicolumn{1}{|l}{x + \alpha } \\ \cline{8-8}
    x^3 & + & \alpha x^2 &   &            &   &                    & \multicolumn{1}{|l}{x^2 + \alpha x + (\alpha ^2 + 1)} \\ \cline{1-3}
        &   & \alpha x^2 & + & x          & + & 1                  & \\
        &   & \alpha x^2 & + & \alpha^2 x      &   &                    & \\ \cline{3-5}
        &   &      &   & (\alpha^2 + 1)x & + & 1                  & \\
        &   &      &   & (\alpha^2 + 1)x & + & \alpha(\alpha^2 + 1)         & \\ \cline{5-7}
        &   &      &   &            &   & \multicolumn{1}{c}{\alpha^3 + \alpha + 1 = 0} &
    \end{array}
    \]
    Buscamos ahora una raíz de $x^2+\alpha x+(\alpha ^2+1)$. Probamos con $\alpha ^2$ (donde usamos que $\alpha^3 + \alpha + 1 = 0$):
    \begin{equation*}
        {\left(\alpha^2\right)}^{2}+ \alpha \alpha^2 + (\alpha^2+1) = \alpha^4 + \alpha^3 + \alpha^2 + 1 = \alpha^4 + \alpha + \alpha^2 = \alpha(\alpha^3 + \alpha + 1) = 0
    \end{equation*}
    Dividimos ahora entre $x+\alpha^2$:
        \[
        \renewcommand{\arraystretch}{1.2}
        \setlength{\arraycolsep}{2pt}
        \begin{array}{rrrrr@{\quad}l}
        x^2 & + & \alpha x               & + & (\alpha^2 + 1)                      & \multicolumn{1}{|l}{x + \alpha^2} \\ \cline{6-6}
        x^2 & + & \alpha^2 x             &   &                                     & \multicolumn{1}{|l}{x + (\alpha^2 + \alpha)} \\ \cline{1-3}
            &   & (\alpha^2 + \alpha)x   & + & (\alpha^2 + 1)                      & \\
            &   & (\alpha^2 + \alpha)x   & + & \alpha^2(\alpha^2 + \alpha)         & \\ \cline{3-5}
            &   &                        &   & \multicolumn{1}{c}{\alpha^4 + \alpha^3 + \alpha^2 + 1 = 0} & 
        \end{array}
        \]
    En definitiva, la factorización de $f^\sigma$ en $\sigma(\bb{Z}_2)[x]$ es:
    \begin{equation*}
        x^3+x+1 = (x+\alpha)(x+\alpha^2)(x+(\alpha^2+\alpha)) \in \sigma(\bb{Z}_2)[x]
    \end{equation*}

    De esta forma, tenemos que $K = \sigma(\bb{Z}_2)(\alpha,\alpha^2,\alpha^2+\alpha)= \sigma(\bb{Z}_2)(\alpha)$ es un cuerpo de descomposición de $f^\sigma$.
\end{ejemplo}

Notemos que, en estos dos ejemplos, hemos visto que, para obtener un cuerpo finito de orden $p^n$, basta con tomar un polinomio irreducible de grado $n$ sobre $\bb{Z}_p$ y construir el cuerpo cociente correspondiente. De hecho, en el Teorema~\ref{teo:clasificacion_cuerpos_finitos} veremos que todos los cuerpos finitos son de esta forma (salvo isomorfismos).\\

Abordamos ahora el problema de notar el cuerpo de descomposición como ``un''. Veremos en el Teorema~\ref{teo:UnicidadCuerpoDescomposicion} que es único salvo isomorfismos.
\begin{lema}\label{lema:extension}
    Sea $F\stackrel{\sigma}{\to}{K}$ un homomorfismo de cuerpos, $p\in F[x]$ irreducible y $\alpha\in K$ una raíz de $p^\sigma$, entonces se tiene que:
    \Func{\sigma_\alpha}{\dfrac{F[x]}{\langle p\rangle}}{\sigma(F)(\alpha)}{g+\langle p\rangle}{g^\sigma(\alpha)}
    es un isomorfismo de cuerpos.
    \begin{proof}
        Podemos tomar:
        \Func{\overline{\sigma_\alpha}}{F[x]}{K}{g}{g^\sigma(\alpha)}
        
        Como $K$ es un cuerpo, este homomorfismo de anillos está bien definido.\\

        Como $p^\sigma(\alpha) = 0$ tenemos que $p\in \ker(\overline{\sigma_\alpha})$, y como $\overline{\sigma_\alpha}$ es un homomorfismo de anillos tenemos que $\ker(\overline{\sigma_\alpha})$ es un ideal, que tiene que ser principal por ser $F[x]$ un Dominio Euclídeo, luego existe $f\in F[x]$ con $\ker(\overline{\sigma_\alpha}) = \langle f \rangle $, de donde ha de existir $h\in F[x]$ con $p = hf$, pero como $p$ es irreducible tiene que ser $h\in \cc{U}(F[x])$, de donde $f = h^{-1}p$, por lo que también $\langle p \rangle = \langle f \rangle =\ker(\overline{\sigma_\alpha})$.\\


        Veamos ahora que $\sigma(F)(\alpha)= Im\overline{\sigma_\alpha}$.
        \begin{description}
            \item[$\supset$)] Sea $f^\sigma(\alpha)\in Im\overline{\sigma_\alpha}$, entonces $\exists f\in F[x]$, $f=\sum\limits_i f_i x^i$; tal que se tiene que $f^\sigma(\alpha) = \overline{\sigma_\alpha}(f) \in Im\overline{\sigma_\alpha}$. Por tanto:
            \begin{equation*}
                f^\sigma(\alpha) = \sum_i f_i^\sigma \alpha^i \in \sigma(F)(\alpha)
            \end{equation*}
            \item[$\subset$)] Sea $a\in F$, entonces $\sigma(a) = a^\sigma(\alpha) = \overline{\sigma_\alpha}(a)$, luego $\sigma(F)\subseteq Im\overline{\sigma_\alpha}$. Además, $\alpha = x^\sigma(\alpha) = \overline{\sigma_\alpha}(x)$, por lo que $\alpha\in Im\overline{\sigma_\alpha}$. Por tanto, $\sigma(F)\cup \{\alpha\}\subseteq Im\overline{\sigma_\alpha}$. Para ver que $\sigma(F)(\alpha)\subseteq Im\overline{\sigma_\alpha}$, basta ver que $Im\overline{\sigma_\alpha}$ es un cuerpo.
            
            Para ello, aplicamos el Primer Teorema de Isomorfía para anillos:
            \begin{equation*}
                \dfrac{F[x]}{\langle p \rangle } = \dfrac{F[x]}{\ker(\overline{\sigma_\alpha})} \cong Im \overline{\sigma_\alpha}
            \end{equation*}
            Como $F[x]/\langle p \rangle$ es un cuerpo (por ser $p$ irreducible), tenemos que $Im \overline{\sigma_\alpha}$ es un cuerpo, y por tanto $\sigma(F)(\alpha)\subseteq Im\overline{\sigma_\alpha}$.
        \end{description}

        Por tanto, $\sigma(F)(\alpha) = Im\overline{\sigma_\alpha}$, luego por el Primer Teorema de Isomorfía para anillos, tenemos que $\sigma_\alpha$ es un isomorfismo de cuerpos.
    \end{proof}
\end{lema}

\begin{definicion}[Extensión de homomorfismos]
    Consideramos dos homomorfismos de cuerpos:
    \begin{figure}[H]
        \centering
        \shorthandoff{""}
        \begin{tikzcd}
        F \arrow[r, "\tau"] \arrow[rd, "\sigma"'] & E                    \\
                                                  & K 
        \end{tikzcd}
        \shorthandon{""}
    \end{figure}
    \noindent
    Diremos que un homomorfismo de cuerpos $\eta:K\to E$ es una $\sigma-$extensión de $\tau$ si:
    \begin{equation*}
        \eta\sigma = \tau
    \end{equation*}
    \noindent
    Y notaremos al conjunto de todas las $\sigma-$extensiones de $\tau$ por:
    \begin{equation*}
        \Ex(\tau,\sigma) = \{\eta:K\to E \text{\ con\ } \eta \text{\ homomorfismo y\ } \eta \sigma = \tau\}
    \end{equation*}
    Notemos que todos estos hacen que el siguiente diagrama sea conmutativo:
    \begin{figure}[H]
        \centering
        \shorthandoff{""}
        \begin{tikzcd}
        F \arrow[r, "\tau"] \arrow[rd, "\sigma"'] & E                    \\
                                                  & K \arrow[u, "\eta"']
        \end{tikzcd}
        \shorthandon{""}
    \end{figure}
    \noindent
    Si $F\leq K$ es una extensión de cuerpos y $\sigma:F\hookrightarrow K$ es la aplicación inclusión, diremos simplemente que $\Ex(\tau,\sigma)$ es el conjunto de las extensiones de $\tau$.
\end{definicion}

Notemos que el motivo de denominar a $\eta$ una ``extensión'' se entiende mejor en el caso de $F\leq K$ y $\sigma:F\hookrightarrow K$, ya que en dicho caso, $\eta$ es una extensión de $\tau$ en el sentido de que $\eta$ es ``la misma aplicación que $\tau$'' pero con un dominio más grande, ya que $\eta$ está definido en $K$ y $\tau$ en $F$.

\begin{prop}[Extensión I]
    Consideremos dos homomorfismos de cuerpos:
    \begin{figure}[H]
        \centering
        \shorthandoff{""}
        \begin{tikzcd}
        F \arrow[r, "\tau"] \arrow[rd, "\sigma"'] & E                    \\
                                                  & K 
        \end{tikzcd}
        \shorthandon{""}
    \end{figure}
    \noindent
    Sean $p\in F[x]$ irreducible y $\alpha\in K$ con $p^{\sigma}(\alpha)=0$. Si denotamos por $\cc{R}$ al conjunto de todas las raíces de $p^\tau$ en $E$ y tenemos que $K = \sigma(F)(\alpha)$, tenemos entonces que la aplicación
    \Func{\psi}{\Ex(\tau,\sigma)}{\cc{R}}{\eta}{\eta(\alpha)}
    es una biyección.
    \begin{proof}
        Veamos en primer lugar que dicha aplicación está bien definida. Para ello, sea $\eta \in \Ex(\tau, \sigma)$:
        \begin{equation*}
            p^\tau(\eta(\alpha)) = p^{\eta \sigma}(\eta (\alpha)) \AstIg \eta(p^\sigma(\alpha)) = \eta(0) = 0
        \end{equation*}
        donde en $(\ast)$ hemos usado que si $p$ es de la forma:
        \begin{equation*}
            p = \sum_i p_i x^i, \qquad p_i \in F
        \end{equation*}

        Entonces, por ser $\eta$ un homomorfismo de cuerpos, tenemos que:
        \begin{equation*}
            p^{\eta\sigma}(\eta(\alpha)) = \sum_i \eta(\sigma(p_i)) \eta(\alpha)^i = \sum_i \eta(\sigma(p_i)\alpha^i) = \eta\left(\sum_i \sigma(p_i) \alpha^i\right) = \eta(p^\sigma(\alpha))
        \end{equation*}
        Esto prueba que\footnote{Notemos que hemos probado además que $\Ex(\tau,\sigma)\neq \emptyset \Longrightarrow \cc{R}\neq \emptyset $.} $\eta(\alpha)\in \cc{R}$. Veamos ahora que la aplicación enunciada es sobreyectiva\footnote{Con lo que tendremos $\cc{R}\neq \emptyset \Longrightarrow \Ex(\tau,\sigma)\neq \emptyset $.}. Para ello, sea $\beta\in \cc{R}$, buscamos una $\sigma-$extensión $\eta$ de $\tau$ de forma que $\eta(\alpha) = \beta$. Usando el Lema~\ref{lema:extension}, obtenemos los isomorfismos:
        \Func{\sigma_\alpha }{\dfrac{F[x]}{\langle p \rangle }}{\sigma(F)(\alpha)=K}{g+\langle p\rangle}{g^\sigma(\alpha)}
        \Func{\tau_\beta }{\dfrac{F[x]}{\langle p \rangle }}{\tau(F)(\beta)\leq E}{g+\langle p\rangle}{g^\tau(\beta)}
        Si tomamos:
        \begin{equation*}
            \eta = i \circ \tau_\beta \circ \sigma_\alpha^{-1}
        \end{equation*}
        donde $i$ es la inclusión $\tau(F)(\beta)\leq E$, observamos que:
        \begin{equation*}
            K\stackrel{\sigma_\alpha^{-1}}{\longrightarrow} \frac{F[x]}{\langle p \rangle } \stackrel{\tau_\beta}{\longrightarrow} \tau(F)(\beta) \stackrel{i}{\longrightarrow} E
        \end{equation*}
        Comprobemos que $\eta\in \Ex(\tau,\sigma)$, ya que si $a\in F$:
        \begin{equation*}
            (\eta \circ \sigma)(a) = (i\circ \tau_\beta\circ \sigma_\alpha^{-1})(\sigma(a)) = (i\circ \tau_\beta)(\sigma_\alpha^{-1}(\sigma(a))) = (i\circ \tau_\beta)(a+\langle p \rangle ) = i(\tau(a)) = \tau(a)
        \end{equation*}
        donde hemos aplicado que tanto $\sigma_\alpha$ como $\tau_\beta$ aplicado sobre constantes son iguales a $\sigma$ y a $\tau$, respectivamente, lo que prueba que $\eta\in \Ex(\tau,\sigma)$. Ahora:
        \begin{equation*}
            \eta(\alpha) = (i\circ \tau_\beta)(\sigma_\alpha^{-1}(\alpha))= (i\circ \tau_\beta)(x+\langle p \rangle ) = i(\beta) = \beta
        \end{equation*}
        Falta probar que la aplicación es inyectiva. Para ello, sean $\eta,\eta' \in \Ex(\tau,\sigma)$ de forma que $\eta(\alpha)=\eta'(\alpha)$, entonces como $\sigma(F)\leq K=\sigma(F)(\alpha)$ con $\alpha$ algebraico sobre $\sigma(F)$, tenemos que $\{1,\alpha,\alpha^2, \ldots\}$ es un sistema de generadores de $\sigma(F)(\alpha)$, por lo que todo elemento de este cuerpo será de la forma:
        \begin{equation*}
            \sum_{i} \sigma(a_i) \alpha^i \in \sigma(F)(\alpha), \qquad a_i \in F
        \end{equation*}
        con lo que:
        \begin{equation*}
            \eta\left(\sum_{i}\sigma(a_i)\alpha^i\right) = \sum_i \eta(\sigma(a_i)) {\eta(\alpha)}^{i} \AstIg \sum_i \eta'(\sigma(a_i)) {\eta'(\alpha)}^{i} = \eta'\left(\sum_{i}\sigma(a_i)\alpha^i\right) 
        \end{equation*}
        donde en $(\ast)$ usamos que $\eta(\alpha) = \eta'(\alpha)$, así como que $\eta, \eta'$ son $\sigma-$extensiones de $\tau$, con lo que $\eta\circ \sigma = \tau = \eta'\circ \sigma$. En definitiva, tenemos que $\eta = \eta'$, al ser $\eta(g) = \eta'(g)$ para todo $g\in K$, lo que nos dice que la aplicación es inyectiva.
    \end{proof}
\end{prop}

\noindent
Obsevemos que en esta última proposición hemos probado además que:
\begin{equation*}
    \cc{R}= \emptyset  \Longleftrightarrow \Ex(\tau,\sigma) =  \emptyset 
\end{equation*}

\begin{lema}\label{lema:extension_compuesta}
    Sean tres homomorfismos entre cuerpos:
    \begin{figure}[H]
        \centering
        \shorthandoff{""}
        \begin{tikzcd}
        F \arrow[r, "\tau"] \arrow[d, "\sigma_1"'] & L   \\
        E_1 \arrow[r, "\sigma_2"']                 & E_2
        \end{tikzcd}
        \shorthandon{""}
    \end{figure}
    Se verifica que:
    \begin{equation*}
        \Ex(\tau,\sigma_2\sigma_1) = \biguplus_{\eta \in \Ex(\tau,\sigma_1)} \Ex(\eta,\sigma_2)
    \end{equation*}
    \begin{proof}
        Por doble inclusión:
        \begin{description}
            \item [$\subseteq )$] Si tomamos $\theta \in \Ex(\tau,\sigma_2\sigma_1)$, tenemos entonces que:
                \begin{equation*}
                    \theta\sigma_2\sigma_1 = \tau
                \end{equation*}
                Por lo que si tomamos $\eta = \theta\sigma_2$, tenemos que:
                \begin{align*}
                    \eta\sigma_1 =\theta\sigma_2\sigma_1 = \tau &\Longrightarrow \eta \in \Ex(\tau,\sigma_1) \\
                    \theta\sigma_2 = \eta &\Longrightarrow \theta \in \Ex(\eta, \sigma_2)
                \end{align*}
            \item [$\supseteq )$] Si $\eta \in \Ex(\tau,\sigma_1)$ y tomamos $\theta\in \Ex(\eta,\sigma_2)$, tendremos entonces que:
                \begin{equation*}
                    \left.\begin{array}{l}
                        \eta\sigma_1 = \tau \\
                        \theta\sigma_2 = \eta
                    \end{array}\right\}\Longrightarrow \theta\sigma_2\sigma_1 = \tau \Longrightarrow \theta\in \Ex(\tau,\sigma_1\sigma_2)
                \end{equation*}
        \end{description}
        \begin{figure}[h]
            \centering
            % https://tikzcd.yichuanshen.de/#N4Igdg9gJgpgziAXAbVABwnAlgFyxMJZABgBpiBdUkANwEMAbAVxiRADEQBfU9TXfIRQBGclVqMWbADLdeIDNjwEiZYePrNWiEAFEA+sLl8lgoqPXVNUnQYBM3cTCgBzeEVAAzAE4QAtkhkIDgQSKISWmwAOlE4dEzGID7+gdQhSHZWktogMdgufnSGINQMdABGMAwACvzKQiDeWC4AFjiJyQGImcGhiADMWZE6ec2F+g6lFVW1pio6Ta3tPF6+XYO9YUM2ubEtMHElIGWVNXVmC81tHWsZaX3h1jkxB3SOXEA
            \begin{tikzcd}
            F \arrow[r, "\tau"] \arrow[d, "\sigma_1"']    & L                        \\
            E_1 \arrow[r, "\sigma_2"'] \arrow[ru, "\eta"] & E_2 \arrow[u, "\theta"']
            \end{tikzcd}
            \caption{Diagrama de la demostración del Lema~\ref{lema:extension_compuesta}.}
            \label{fig:extension_compuesta}
        \end{figure}

        La Figura~\ref{fig:extension_compuesta} ilustra la demostración. Hemos probado que
        \begin{equation*}
            \Ex(\tau,\sigma_2\sigma_1) = \bigcup_{\eta \in \Ex(\tau,\sigma_1)} \Ex(\eta,\sigma_2)
        \end{equation*}
        Ahora, si $\eta,\eta'\in \Ex(\tau,\sigma_1)$ y tenemos que:
        \begin{equation*}
            \theta \in  \Ex(\eta,\sigma_2)\cap \Ex(\eta',\sigma_2) \Longrightarrow \left\{\begin{array}{l}
                \theta\sigma_2 = \eta \\
                \theta \sigma_2 = \eta'
            \end{array}\right. \Longrightarrow \eta = \eta'
        \end{equation*}
        por lo que la unión es disjunta.
    \end{proof}
\end{lema}

\begin{prop}[Extensión II]\label{prop:extension_2}
    Sean dos homomorfismos de cuerpos:
    \begin{figure}[H]
        \centering
        \shorthandoff{""}
        \begin{tikzcd}
        F \arrow[r, "\tau"] \arrow[rd, "\sigma"'] & E \\
                                                  & K
        \end{tikzcd}
        \shorthandon{""}
    \end{figure}
    Si $[K:\sigma(F)]<\infty$, entonces $|\Ex(\tau,\sigma)| \leq [K:\sigma(F)]$.
    \begin{proof}
        Por inducción sobre $n = [K:\sigma(F)]$ usando el segundo principio de inducción:
        \begin{itemize}
            \item \underline{Si $n=1$}, entonces $\sigma(F) = K$, por lo que $\sigma$ es un isomorfismo, con lo que $\Ex(\tau,\sigma) = \{\tau \sigma^{-1}\}$, ya que si $\eta \in \Ex(\tau, \sigma)$, entonces:
                \begin{equation*}
                    \eta\sigma = \tau \Longrightarrow \eta = \tau\sigma^{-1}
                \end{equation*}
                En definitiva, $1 = |\Ex(\tau,\sigma)| \leq [K:\sigma(F)] = 1$.
            \item \underline{Supuesto que $n>1$ y la hipótesis de inducción}, como $[K:\sigma(F)] = n > 1$, tenemos que existe $\alpha\in K$ de forma que $[\sigma(F)(\alpha):\sigma(F)] > 1$, con lo que el Lema de la Torre nos dice que $[K:\sigma(F)(\alpha)] < n$. 

                Sea ahora $\iota:\sigma(F)(\alpha)\to K$ la inclusión en $K$, podemos tomar:
                \begin{equation*}
                    \sigma = \iota \circ \sigma'
                \end{equation*}
                con $\sigma':F\to \sigma(F)(\alpha)$ la restricción en codominio (o correstricción) de $\sigma$. Nos encontramos en la siguiente situación:
                \begin{figure}[H]
                    \centering
                    \shorthandoff{""}
                    \begin{tikzcd}
                        F \arrow[r, "\tau"] \arrow[rd, "\sigma"] \arrow[d, "\sigma'"'] & E \\
                        \sigma(F)(\alpha) \arrow[r, "\iota"]                           & K
                    \end{tikzcd}
                    \shorthandon{""}
                \end{figure}
                \noindent
                Aplicando el Lema~\ref{lema:extension_compuesta}, obtenemos:
                \begin{equation*}
                    \Ex(\tau,\sigma) = \biguplus_{\eta \in \Ex(\tau,\sigma')}\Ex(\eta, \iota)
                \end{equation*}
                Con lo que:
                \begin{equation*}
                    |\Ex(\tau,\sigma)| = \sum_{\eta \in  \Ex(\tau, \sigma')}|\Ex(\eta,\iota)|
                \end{equation*}

                Buscamos ahora una cota superior para ambos cardinales.
                \begin{description}
                    \item [Cota superior para $|\Ex(\eta,\iota)|$:] Sea se tiene que $\eta\in \Ex(\tau,\sigma')$. Entonces $\eta:\sigma(F)(\alpha)\to E$ es un homomorfismo de cuerpos, y tenemos la siguiente situación:
                    \begin{figure}[H]
                        \centering
                        \shorthandoff{""}
                        \begin{tikzcd}
                        \sigma(F)(\alpha) \arrow[r, "\eta"] \arrow[rd, "\iota"'] & E \\
                                                                  & K
                        \end{tikzcd}
                        \shorthandon{""}
                    \end{figure}
                    Como $[K:\sigma(F)(\alpha)] < n$, aplicando la hipótesis de inducción, tenemos que:
                        \begin{equation*}
                            |\Ex(\eta,\iota)| \leq [K:\sigma(F)(\alpha)]
                        \end{equation*}
                    \item [Cota superior para $|\Ex(\tau,\sigma')|$:]
                    
                    Como $[K:\sigma(F)] = n$ finita, es algebraica, y por tanto $\alpha$ es algebraico sobre $\sigma(F)$, luego podemos considerar $\Irr(\alpha,\sigma(F))$. Considerando la biyección $F\to \sigma(F)$, existe un polinomio irreducible $p\in F[x]$ de forma que $p^\sigma = \Irr(\alpha,\sigma(F))$. La Proposición de Extensión nos dice que, si $\cc{R}_\tau$ es el número de raíces de $p^\tau$ en $E$, entonces:
                    \begin{equation*}
                        |\Ex(\tau, \sigma')| = |\cc{R}_\tau|
                    \end{equation*}

                    Como el número de raíces de un polinomio no puede ser mayor que su grado, y $\deg p^\tau = \deg p^\sigma = [\sigma(F)(\alpha):\sigma(F)]$, tenemos que:
                    \begin{equation*}
                        |\Ex(\tau, \sigma')| = |\cc{R}_\tau| \leq \deg p^\tau = [\sigma(F)(\alpha):\sigma(F)]
                    \end{equation*}
                \end{description}

                Uniendo ambas cotas, obtenemos:
                \begin{align*}
                    |\Ex(\tau,\sigma)| &= \sum_{\eta \in  \Ex(\tau, \sigma')}|\Ex(\eta,\iota)| \leq \sum_{\eta \in  \Ex(\tau, \sigma')}[K:\sigma(F)(\alpha)] \\
                    &= |\Ex(\tau, \sigma')|[K:\sigma(F)(\alpha)] \leq [\sigma(F)(\alpha):\sigma(F)][K:\sigma(F)(\alpha)]
                \end{align*}


                Aplicando el Lema de la Torre, obtenemos:
                \begin{equation*}
                    |\Ex(\tau,\sigma)|  \leq [\sigma(F)(\alpha):\sigma(F)][K:\sigma(F)(\alpha)] = [K:\sigma(F)]
                \end{equation*}\qedhere
        \end{itemize}
    \end{proof}
\end{prop}

\begin{ejercicio} \label{ej:homomorfismo_primo}
    Si $\Pi$ es el cuerpo primo de un cuerpo $K$, entonces el único homomorfismo de cuerpos $\sigma:\Pi\to K$ es la inclusión.\\ 

    \noindent
    Sea $\sigma:\Pi\to K$ un homomorfismo de cuerpos, sea $\iota:\Pi\to K$ el homomorfismo inclusión, vemos que:
    \begin{itemize}
        \item $\sigma(0) = 0 = \iota(0)$.
        \item $\sigma(1) = 1 = \iota(1)$.
        \item Si $n\in \mathbb{N}$, tenemos que:
            \begin{equation*}
                \sigma\left(\sum_{k=1}^{n}1\right) = \sum_{k=1}^{n}\sigma(1) = \sum_{k=1}^{n}1 = \sum_{k=1}^{n}\iota(1) = \iota\left(\sum_{k=1}^{n}1\right)
            \end{equation*}
    \end{itemize}
    Distinguimos casos según la Proposición~\ref{prop:subcuerpo_primo}:
    \begin{description}
        \item [Si $\car(K) > 0$.] Tendremos entonces que existe un isomorfismo $\Phi:\mathbb{Z}_p\to \Pi$ para $p=\car(K)$. Si $a\in \Pi$, tenemos que existe $b\in \mathbb{Z}_p$ de forma que:
            \begin{equation*}
                a = \Phi(b) = \Phi\left(\sum_{k=1}^{b}1\right) = \sum_{k=1}^{b} \Phi (1) = \sum_{k=1}^{b} 1
            \end{equation*}
            por lo que $\sigma(a) = \iota(a)$, para todo $a\in \Pi$, luego $\sigma=\iota$.
        \item [Si $\car(K)=0$.] Tendremos entonces que existe un isomorfismo $\Phi:\mathbb{Q}\to \Pi$. Si $a\in \Pi$, tenemos que existen $z\in \mathbb{Z}, n\in \mathbb{N}\setminus\{0\}$ de forma que:
            \begin{align*}
                a &= \Phi\left(\frac{z}{n}\right) = \Phi(z){(\Phi(n))}^{-1} = \Phi\left(\sgn(z)\sum_{k=1}^{|z|}1\right){\left(\Phi\left(\sum_{k=1}^{n}1\right)\right)}^{-1} \\ 
                  &= \sgn(z)\left(\sum_{k=1}^{|z|}\Phi(1)\right){\left(\sum_{k=1}^{n}\Phi(1)\right)}^{-1}= \sgn(z)\left(\sum_{k=1}^{|z|}1\right){\left(\sum_{k=1}^{n}1\right)}^{-1}
            \end{align*}
            por lo que $\sigma(a) = \iota(a)$, para todo $a\in \Pi$, de donde $\sigma=\iota$.
    \end{description}
\end{ejercicio}

\noindent
Mostramos a continuación un ejemplo básico de la proposición de extensión.

\begin{ejemplo}
    ¿Cuántos homomorfismos de cuerpos hay de $\mathbb{Q}\left(\sqrt[3]{2}\right)$ en $\mathbb{C}$, y cuáles son?\\

    \noindent
    Para responder a esta pregunta trataremos de reformularla en una que podamos responder usando la teoría de extensiones de homomorfismos. Sea $\eta:\mathbb{Q}\left(\sqrt[3]{2}\right)\to\mathbb{C}$ un homomorfismo de cuerpos. Si restringimos $\eta$ al cuerpo primo de $\mathbb{Q}\left(\sqrt[3]{2}\right)$, que es $\mathbb{Q}$, el Ejercicio~\ref{ej:homomorfismo_primo} nos dice que entonces $\eta\big|_{\mathbb{Q}}$ coincide con la aplicación inclusión $\tau:\mathbb{Q}\hookrightarrow \mathbb{C}$, es decir:
    \begin{equation*}
        \eta  \circ \iota = \tau
    \end{equation*}
    Si tomamos $\sigma=\iota:\mathbb{Q}\to\mathbb{Q}\left(\sqrt[3]{2}\right)$ la aplicación inclusión, estudiar cuántos homomorfismos de cuerpos hay de $\mathbb{Q}\left(\sqrt[3]{2}\right)$ en $\mathbb{C}$ es equivalente a estudiar cuántas $\sigma-$extensiones de $\tau$ hay.
    \begin{figure}[H]
        \centering
        \shorthandoff{""}
        \begin{tikzcd}
            \mathbb{Q} \arrow[rd, "\sigma"] \arrow[r, "\tau"] & \mathbb{C}                                   \\
                                                             & {\mathbb{Q}\left(\sqrt[3]{2}\right)} \arrow[u, "\eta"']
        \end{tikzcd}
        \shorthandon{""}
    \end{figure}
    \noindent
    Por tanto, el problema se reduce a estudiar los elementos del conjnto $\Ex(\tau,\sigma)$. Sabemos que:
    \begin{equation*}
        \Irr\left(\sqrt[3]{2},\mathbb{Q}\right) = x^3-2
    \end{equation*}
    Cuyas raíces en $\mathbb{C}$ son:
    \begin{equation*}
        \cc{R} = \left\{\sqrt[3]{2}, w\sqrt[3]{2}, w^2\sqrt[3]{2}\right\}
    \end{equation*}
    donde $w$ es una raíz cúbica primitiva de la unidad. La Proposición de extensión nos dice entonces que existen exactamente tres homomorfismos de $\mathbb{Q}\left(\sqrt[3]{2}\right)$ en $\mathbb{C}$. Les damos nombre a cada uno de ellos:
    \begin{equation*}
        \Ex(\tau,\sigma) = \{\eta_0,\eta_1,\eta_2\}
    \end{equation*}
    donde $\eta_i$ está determinado (según la proposición de extensión) por:
    \begin{equation*}
        \eta_j\left(\sqrt[3]{2}\right) = w^j \sqrt[3]{2}, \qquad \forall j \in \{0,1,2\}
    \end{equation*}

    Notemos que cada uno de los $\eta_j$ es un homomorfismo definido sobre $\mathbb{Q}\left(\sqrt[3]{2}\right)$ con base $\left\{1,\sqrt[3]{2}, \left(\sqrt[3]{2}\right)^2\right\}$. Como los homomorfismos respetan sumas, productos y llevan la unidad en la unidad, basta definirlos sobre $\sqrt[3]{2}$, ya que:
    \begin{equation*}
        \eta_j\left(a + b\sqrt[3]{2} + c\left(\sqrt[3]{2}\right)^2\right) = a + b\eta_j\left(\sqrt[3]{2}\right) + c\eta_j\left(\sqrt[3]{2}\right)^2
    \end{equation*}



    Como ejemplo de esto último, observemos que podemos calcular:
    \begin{equation*}
        \eta_2\left(\frac{{\sqrt[3]{2}+\left(\sqrt[3]{2}\right)}^{2}}{27}\right) = \dfrac{\eta_2(\sqrt[3]{2}) + {\left(\eta_2\left(\sqrt[3]{2}\right)\right)}^{2}}{\eta_2(27)} = \dfrac{w^2 \sqrt[3]{2} + {\left(w^2\sqrt[3]{2}\right)}^{2}}{27} 
        = \dfrac{w^2 \sqrt[3]{2} + w\sqrt[3]{4}}{27} 
    \end{equation*} 
\end{ejemplo}

\begin{prop}[Extensión III]\label{prop:extension_3}
    Sean dos homomorfismos de cuerpos:
    \begin{figure}[H]
        \centering
        \shorthandoff{""}
        \begin{tikzcd}
        F \arrow[r, "\tau"] \arrow[rd, "\sigma"'] & E \\
                                                  & K
        \end{tikzcd}
        \shorthandon{""}
    \end{figure}
    \noindent
    con $\sigma$ un cuerpo de descomposición de $f\in F[x]$. Si $f^\tau$ se descompone como producto de polinomios lineales en $E[x]$, entonces $\Ex(\tau,\sigma)$ es no vacío. Además, si $f^\sigma$ tiene $\deg f$ raíces distintas, entonces:
    \begin{equation*}
        |\Ex(\tau, \sigma)| = [K:\sigma(F)]
    \end{equation*}
    \begin{proof}
        Como $\sigma$ es un cuerpo de descomposición de $f$, tenemos que $K = \sigma(F)(\alpha_1,\ldots,\alpha_r)$, con $\alpha_i$ raíces de $f^\sigma$. Consideramos la extensión $\sigma(F)\leq K$:
        \begin{itemize}
            \item Esta es finitamente generada, ya que $K = \sigma(F)(\alpha_1,\ldots,\alpha_r)$.
            \item Esta es algebraica, ya que cada $\alpha_i$ es algebraico sobre $\sigma(F)$, al ser raíz de $f^\sigma$.
        \end{itemize}

        Por tanto, $[K:\sigma(F)]<\infty$, y podemos demostrar la proposición por inducción sobre $n=[K:\sigma(F)]$. La idea es similar a la de la Proposición~\ref{prop:extension_2}:
        \begin{itemize}
            \item Para $n = 1$, tenemos que $K = \sigma(F)$, con lo que $\sigma$ es un isomorfismo y tendremos por tanto que $\Ex(\tau,\sigma) = \{\tau\sigma^{-1}\}$.
            \item Supuesto que $n>1$ y la hipótesis de inducción, tenemos que $f$ tiene un factor irreducible $p\in F[x]$ de grado mayor o igual 1. Tomamos una raíz $\alpha\in K$ de $p^\sigma$, de donde $[K:\sigma(F)(\alpha)]<n$. Si consideramos $\sigma':F\to \sigma(F)(\alpha)$ y la inclusión $\iota:\sigma(F)(\alpha)\to K$, tenemos que:
                \begin{equation*}
                    \Ex(\tau,\sigma) = \biguplus_{\eta\in \Ex(\tau,\sigma')}\Ex(\eta,\iota)
                \end{equation*}
                con lo que:
                \begin{equation*}
                    |\Ex(\tau,\sigma)| = \sum_{\eta \in  \Ex(\tau, \sigma')}|\Ex(\eta,\iota)|
                \end{equation*}
                de la proposición de extensión deducimos que $\Ex(\tau,\sigma')$ tiene tantos elementos como raíces de $p^\tau$ hay en $E$. Sin embargo, como $f^\tau$ se factoriza como producto de polinomios de grado 1 en $E[x]$ y $p^\tau$ es un factor de $f^\tau$, en particular $\Ex(\tau,\sigma')\neq \emptyset $, lo que nos permite tomar $\eta \in \Ex(\tau,\sigma')$, y por hipótesis de inducción obtenemos que $\Ex(\eta, \iota)$ es no vacío, con lo que tampoco puede serlo $\Ex(\tau,\sigma)$.

                Además, si $f^\sigma$ tiene $\deg f$ raíces distintas, entonces $p^\sigma$ tiene $\deg p$ raíces distintas, de donde:
                \begin{equation*}
                    |\Ex(\tau,\sigma')| = \cc{R}(p^\sigma) = \deg p^{\sigma} = [\sigma(F)(\alpha):\sigma(F)]
                \end{equation*}
                Por hipótesis de inducción (como $[K:\sigma(F)(\alpha)]<n$), para cada $\eta \in \Ex(\tau,\sigma')$ tenemos que $|\Ex(\eta, \iota)| = [K:\sigma(F)(\alpha)]$, con lo que:
                \begin{equation*}
                    |\Ex(\tau,\sigma)| = \sum_{\eta \in  \Ex(\tau, \sigma')}|\Ex(\eta,\iota)| = [K:\sigma(F)(\alpha)][\sigma(F)(\alpha):\sigma(F)] = [K:\sigma(F)]
                \end{equation*}
        \end{itemize}
    \end{proof}
\end{prop}

\begin{ejemplo}\label{ejemplo:extension_homomorfismos}
    (Continuación del ejemplo anterior)

    \noindent
    Sea $K = \mathbb{Q}\left(\sqrt[3]{2},w\right)$ con $w$ una raíz cúbica primitiva de la unidad, si queremos calcular todos los homomorfismos de $K$ en $\mathbb{C}$, lo que haremos será considerar las respectivas aplicaciones de inclusión $\tau,\sigma_1$ y $\sigma_2$, con lo que tenemos:
    \begin{figure}[H]
        \centering
        \shorthandoff{""}
        \begin{tikzcd}
        \mathbb{Q} \arrow[d, "\sigma_1"'] \arrow[r, "\tau"]                   & \mathbb{C}           \\
        {\mathbb{Q}(\sqrt[3]{2})} \arrow[ru, "\eta_j"] \arrow[r, "\sigma_2"'] & K \arrow[u, "\eta"']
        \end{tikzcd}
        \shorthandon{""}
    \end{figure}
    \noindent
    Y queremos calcular $\Ex(\tau,\sigma_2\sigma_1)$. Para ello, trataremos de usar las aplicaciones $\eta_j$ que ya conocemos, que cumplían:
    \begin{equation*}
        \eta_j\left(\sqrt[3]{2}\right) = w^j \sqrt[3]{2} \qquad \forall j \in \{0,1,2\}
    \end{equation*}
    Calcularemos para cada $j$ todas las $\sigma_2-$extensiones de $\eta_j$, ya que:
    \begin{equation*}
        \Ex(\tau,\sigma_2\sigma_1) = \biguplus_{\eta_j \in  \Ex(\tau,\sigma_1)}\Ex(\eta_j, \sigma_2) = \Ex(\eta_0,\sigma_2) \cup \Ex(\eta_1, \sigma_2) \cup \Ex(\eta_2,\sigma_2)
    \end{equation*}
    Para ello, necesitamos calcular el polinomio irreducible de $w$ sobre $\mathbb{Q}\left(\sqrt[3]{2}\right)$ y calcular sus raíces en $\mathbb{C}$, algo que se hizo en el Ejercicio~\ref{ej:Irr(w,Q(sqrt[3]{2}))}:
    \begin{equation*}
        \Irr\left(w,\mathbb{Q}\left(\sqrt[3]{2}\right)\right) = x^2+x+1 \qquad \text{con raíces\ } w, w^2
    \end{equation*}
    Por tanto, tendremos dos $\sigma_2$ extensiones de $\eta_j$ para cada $j \in \{0,1,2\}$:
    \begin{equation*}
        \eta_{j,k}(w) = w^k \qquad k \in \{1,2\}
    \end{equation*}

    En resumen:
    \begin{equation*}
        \Ex(\tau,\sigma_2\sigma_1) = \{\eta_{j,k} : j\in \{0,1,2\}, k\in \{1,2\}\}
    \end{equation*}
    determinadas por
    \begin{equation*}
        \eta_{j,k}\left(\sqrt[3]{2}\right) = w^j\sqrt[3]{2}, \qquad \eta_{j,k}(w)= w^k
    \end{equation*}
    Sabíamos que teníamos que obtener 6 extensiones, puesto que $K$ es cuerpo de descomposición de $x^3-2$, con todas sus raíces distintas y que se descompone como producto de polinomios lineales en $\mathbb{C}$, y $[K:\mathbb{Q}] = 6$ como se vió en el ejemplo de la página~\pageref{ej:Ordenwsqrt3_2}. Observemos, por ejemplo, que:
    \begin{equation*}
        \eta_{j,k}\left(w\sqrt[3]{2}\right) = \eta_{j,k}(w)\eta_{j,k}\left(\sqrt[3]{2}\right) = w^k w^j\sqrt[3]{2} = w^{k+j}\sqrt[3]{2}
    \end{equation*}

    Por último, notemos que este ejemplo es válido si, en lugar de $\bb{C}$, consideramos cualquier cuerpo $E$ tal que $K\leq E$.
\end{ejemplo}

\begin{ejercicio}\label{ej:F-lineal}
    Sean $F\stackrel{\tau}{\to}E\stackrel{\rho}{\to}E$ homomorfismos de cuerpos. Sabemos que $E$ es un $\tau(F)-$espacio vectorial. Se verifica que: 
    \begin{equation*}
        \rho \text{\ es\ } \tau(F)-\text{lineal} \Longleftrightarrow \rho\tau = \tau
    \end{equation*}

    \begin{description}
        \item [$\Longleftarrow )$] Sea $y\in \tau(F)$ y $z\in E$, tenemos que existe $x\in F$ de forma que $\tau(x)=y$, con lo que:
            \begin{align*}
                \rho(y\cdot z) = \rho(\tau(x)\cdot z) = \rho(\tau(x))\cdot \rho(z) = \tau(x)\cdot \rho(z) = y\cdot \rho(z)
            \end{align*}
        \item [$\Longrightarrow )$] Supuesto que $\rho$ es $\tau(F)-$lineal, tenemos que:
            \begin{equation*} 
                \rho(\tau(x)) = \rho(\tau(x)\cdot 1)= \tau(x)\cdot \rho(1) = \tau(x)\cdot 1 = \tau(x) \qquad \forall x\in E
            \end{equation*}
    \end{description}
\end{ejercicio}

Llegamos en este punto a un resultado importante, que nos explica en qué sentido el cuerpo de descomposición de un polinomio es único.
\begin{teo}[Unicidad del cuerpo de descomposición]\label{teo:UnicidadCuerpoDescomposicion} \ \\
    Sean $\tau:F\to E$ y $\tau':F\to E'$ dos cuerpos de descomposición de $f\in F[x]$. Entonces, existe un isomorfismo de cuerpos $\eta:E\to E'$ tal que $\eta \tau = \tau'$.
    \begin{proof}
        La Proposición~\ref{prop:extension_3} nos dice que como $f^\tau$ y $f^{\tau'}$ se descomponen como producto de polinomios lineales en $E[x]$ y $E'[x]$ de forma respectiva, entonces $\Ex(\tau,\tau')$ y $\Ex(\tau',\tau)$ son no vacíos, con lo que existen $\eta:E\to E'$ y $\eta':E'\to E$ tales que 
        \begin{equation*}
            \eta'\tau' = \tau\qquad \eta\tau = \tau'
        \end{equation*}

        Si observamos que $\eta\eta'\tau' = \tau'$, el Ejercicio~\ref{ej:F-lineal} nos dice que $\eta\eta'$ es $\tau'(F)-$lineal. Ahora, como $E'$ es de dimensión finita sobre $\tau'(F)$ por ser $E'$ cuerpo de descomposición de $f^{\tau'}$; y como tenemos que $\eta\eta':E'\to E'$ es inyectiva por la fórmula de las dimensiones, tenemos que:
        \begin{equation*}
            \dim_{\tau'(F)}E' = \dim_{\tau'(F)} Im(\eta\eta') + \dim_{\tau'(F)} \ker(\eta\eta') = \dim_{\tau'(F)} Im(\eta\eta')
        \end{equation*}

        Por tanto, $Im(\eta\eta') = E'$ y $\eta\eta'$ es sobreyectiva, luego $\eta$ es sobreyectiva. Como era un homomorfismo de cuerpos, concluimos que $\eta$ es biyectiva, con lo que $\eta$ es un isomorfismo.
    \end{proof}
\end{teo}

\begin{ejercicio}
    Sea $\sigma:F\to E$ un homomorfismo de cuerpos tal que la extensión $\sigma(F)\leq E$ es finita. Demostrar que existe un polinomio $f\in F[x]$ y un homomorfismo de cuerpos $\tau:E\to K$ tal que $\tau\sigma:F\to K$ es cuerpo de descomposición de $f$.\\

    \noindent
    Como la extensión $\sigma(F)\leq E$ es finita, sabemos entonces que es algebraica y finitamente generada, con lo que existen $\alpha_1, \ldots, \alpha_n\in E$ algebraicos sobre $\sigma(F)$ de forma que $E = \sigma(F)(\alpha_1, \ldots, \alpha_n)$. Obtenemos para todo $i \in \{1,\ldots,n\}$:
    \begin{equation*}
        g_i = \Irr(\alpha_i, \sigma(F))\in \sigma(F)[x]
    \end{equation*}
    con lo que $g_i(\alpha_i)=0\quad \forall i \in \{1,\ldots,n\}$. Como $\sigma:F\to\sigma(F)$ es un isomorfismo, para cada $g_i$ existe un único polinomio $f_i \in F[x]$ de forma que $f_i^\sigma=g_i$. Consideramos:
    \begin{equation*}
        f = \prod_{i=1}^{n} f_i \Longrightarrow f^\sigma =\prod_{i=1}^{n}f_i^\sigma= \prod_{i=1}^{n}g_i \in \sigma(F)[x]
    \end{equation*}
    Por la Proposición~\ref{prop:existe_cdd}, sabemos que podemos encontrar $\theta:\sigma(F)\to K$ cuerpo de descomposición de $f^{\sigma}$. Trataremos ahora de extender $\theta$ a $E$. Para ello, si observamos que:
    \begin{figure}[H]
        \centering
        \shorthandoff{""}
        \begin{tikzcd}
            \sigma(F) \arrow[r, "\theta"] \arrow[rd, "\iota"] & K                   \\
                                                            & \sigma(F)(\alpha_1)
        \end{tikzcd}
        \shorthandon{""}
    \end{figure}
    \noindent
    y recordamos que $g_1\in \sigma(F)[x]$ es irreducible en $\sigma(F)$, la proposición de extensión nos dice que existe existe $\eta_1 \in \Ex(\theta,\iota)$ de forma que $\eta_1(\alpha_1)$ es una raíz de $g_1^\theta$ (y por tanto de ${(f^\sigma)}^{\theta}$) en $K$. Supuesto ahora que:
    \begin{figure}[H]
        \centering
        \shorthandoff{""}
        \begin{tikzcd}
            \sigma(F)(\alpha_1, \ldots,\alpha_k) \arrow[r, "\eta_k"] \arrow[rd, "\iota"] & K                   \\
                                                                                         & \sigma(F)(\alpha_1, \ldots, \alpha_{k+1})
        \end{tikzcd}
        \shorthandon{""}
    \end{figure}
    Si tomamos $\Irr(\alpha_{k+1},\sigma(F)(\alpha_1, \ldots, \alpha_k))$ (divisor de $g_{k+1}$), la proposición de extensión nos garantiza la existencia de $\eta_{k+1}\in \Ex(\eta_k,\iota)$ de forma que $\eta_{k+1}(\alpha_{k+1})$ es una raíz de $g_{k+1}^{\eta_k}$. Tomando ahora $\tau = \eta_{n}$, tenemos $\tau:\sigma(F)(\alpha_1, \ldots, \alpha_n) = E\to K$ de forma que ${(f^{\sigma})}^{\tau}$ se descompone como producto de polinomios lineales en $K[x]$, y si $\cc{R}$ es el conjunto de raíces de ${(f^{\sigma})}^{\tau}$, $K = \sigma(F)(\cc{R})$, con lo que $K$ es cuerpo de descomposición de ${(f^{\sigma})}^{\tau}$.
\end{ejercicio}

\section{Clasificación de los cuerpos finitos} 
\begin{prop}\label{prop:cuerpo_descomposicion}
    Sea $F$ un cuerpo finito de cardinal\footnote{Sabemos que es así por el Ejercicio~\ref{ej:cardinal_cuerpo}.} $q=p^n$  donde $p=\car(F)$, entonces $F$ es cuerpo de descomposición de $x^q-x\in \bb{Z}_p[x]$.
    \begin{proof}
        Llamamos $f=x^q-x$ y consideramos el grupo $F^{\times} = F\setminus\{0\}$, que tiene $q-1$ elementos. Por el Teorema de Lagrange para grupos tenemos que todo $\alpha\in F^{\times}$ satisface que $\alpha^{q-1}=1$, de donde $\alpha^q = \alpha$. Para $0$ es trivial, con lo que:
        \begin{equation*}
            \alpha^q = \alpha \qquad \forall \alpha\in F
        \end{equation*}
        es decir, todo elemento de $F$ es raíz de $x^q-x$. Como $q=p^n=0\in \bb{Z}_p$, tenemos que su polinomio derivado es:
        \begin{equation*}
            f' = qx^{q-1}-1 = -1\in \bb{Z}_p[x]
        \end{equation*}
        
        Tenemos entonces que $x^q-x$ tiene exactamente $q$ raíces distintas, que deben ser todos aquellos elementos de $F$. Esto prueba que $f$ se descompone como producto de polinomios lineales en $F[x]$:
        \begin{equation*}
            f = \prod_{\alpha\in F} (x-\alpha) \in F[x]
        \end{equation*}
        
        Falta por ver que $F=\bb{Z}_p(\cc{R})$ con $\cc{R}$ el conjunto de raíces de $f$. Como $\cc{R} = F$, tenemos que $\bb{Z}_p(\cc{R}) = \bb{Z}_p(F)$. Por otro lado, como $\mathbb{Z}_p\leq F$, se tiene $\bb{Z}_p(F) = F$.
    \end{proof}
\end{prop}

\begin{ejercicio}\label{ej:cuerpo_caracteristica}
    Sean $a,b\in F$ con $F$ un cuerpo de característica $p>0$. Si $q=p^n$, comprobar que ${(a-b)}^{q} = a^q - b^q$.\\

    \noindent
    Veamos en primer lugar que:
    \begin{equation*}
        {(a-b)}^{p} = \sum_{k=0}^{p} \binom{p}{k} a^{p-k}{(-b)}^k = a^p - b^p + \sum_{k=1}^{p-1} \dfrac{p!}{k!(p-k)!}a^{p-k}{(-b)}^{k} \AstIg a^p - b^p
    \end{equation*}
    donde en $(\ast)$ usamos que para $1<k<p-1$ tenemos que $\binom{p}{k}$ es múltiplo de $p$. Observemos ahora que:
    \begin{equation*}
        {(a-b)}^{p^2} = {({(a-b)}^{p})}^{p} = {(a^p-b^p)}^{p} = a^{p^2} - b^{p^2}
    \end{equation*}
    Y por un procedimiento inductivo se termina probando que ${(a-b)}^{q} = a^q - b^q$.
\end{ejercicio}

\begin{teo}[Clasificación de cuerpos finitos]\label{teo:clasificacion_cuerpos_finitos}
    Para cada número primo $p$ y para cada $n\in \mathbb{N}\setminus\{0\}$ existe un único, salvo isomorfismos, cuerpo de cardinal $p^n$. Además, estos son los únicos cuerpos finitos.
    \begin{proof}
        Sea $q=p^n$, tomamos como $F$ un cuerpo de descomposición del polinomio $f=x^q-x\in \bb{Z}_p[x]$, que sabemos que existe por la Proposición~\ref{prop:existe_cdd}. Hemos de ver que $|F|=q$. Sea $S$ el conjunto de las raíces de $f$ en $F$, veamos que $S$ es un subcuerpo de $F$:
        \begin{itemize}
            \item $1\in S$.
            \item Si $a,b\in S$:
                \begin{equation*}
                    \left.\begin{array}{l}
                        a^q - a = 0 \\
                        b^q - b = 0 
                    \end{array}\right\} \Longrightarrow a^qb^q = ab \Longrightarrow {(ab)}^{q}-ab = 0
                \end{equation*}
                con lo que $ab\in S$, y vemos ahora que:
                \begin{equation*}
                    {(a-b)}^{q}-(a-b) \AstIg a^q - b^q - (a-b) = a-b-(a-b) = 0
                \end{equation*}
                donde en $(\ast)$ usamos el Ejercicio~\ref{ej:cuerpo_caracteristica}, con lo que también $a-b\in S$.
            \item Ahora, si $a\in S\setminus\{0\}$, tenemos que:
                \begin{equation*}
                    {(a^{-1})}^{q} = a^{-q} = {(a^{q})}^{-1} = a^{-1} \Longrightarrow {(a^{-1})}^{q}-a^{-1} = 0
                \end{equation*}
                por lo que $a^{-1}\in S$.
        \end{itemize}
        Como $\mathbb{Z}_p\leq F$ es el cuerpo primo y $S\leq F$, ha de ser $\mathbb{Z}_p\leq S$. Finalmente, como $F$ es un cuerpo de descomposición de $f$, ha de ser $F=\mathbb{Z}_p(S) =S$. Además, como el polinomio derivado no comparte raíces con $f$, tenemos que $|F| = q$.\\

        \noindent
    Ahora, si tenemos dos cuerpos del mismo cardinal $q$, la Proposición~\ref{prop:cuerpo_descomposicion} nos dice que ambos cuerpos son cuerpos de descomposición de $x^q-x\in \bb{Z}_p[x]$, y aplicando el Teorema de unicidad del cuerpo de descomposición, tenemos que son isomorfos.\\

    \noindent
    Sea ahora $F$ cualquier cuerpo finito, tenemos por el Ejercicio~\ref{ej:cardinal_cuerpo} que este tiene cardinal $p^n$, por lo que tenemos el resultado por lo que acabamos de probar.
    \end{proof}
\end{teo}

\begin{notacion}
    Si $F$ es un cuerpo de $q=p^n$ elementos, lo notaremos por $\bb{F}_q$, y hablaremos ``del'' cuerpo de $q$ elementos. Como todos los cuerpos de $q$ elementos son isomorfos entre sí, usaremos $\bb{F}_q$ como una etiqueta que hace referencia a cualquier cuerpo de $q$ elementos. Estos cuerpos reciben el nombre de \underline{cuerpos de Galois}.
\end{notacion}

\section{El grupo de automorfismos de una extensión}
\begin{definicion}[Grupo de automorfismos de un cuerpo]
    Sea $K$ un cuerpo, consideramos el conjunto de todos los automorfismos de $K$:
    \begin{equation*}
        \Aut(K) = \{\sigma:K\to K \text{\ homomorfismo de cuerpos biyectivo}\}
    \end{equation*}
    Se verifica que $\Aut(K)$ es un grupo con la operación composición de aplicaciones, que recibe el nombre de \underline{grupo de automorfismos de $K$}. \\

    \noindent
    Si $F\leq K$ es una extensión de cuerpos, consideraremos también:
    \begin{equation*}
        \Aut_F(K) = \{\sigma\in \Aut(K) : \sigma \text{\ es\ } F-\text{lineal}\}
    \end{equation*}
    y se verifica que $\Aut_F(K)$ es un subgrupo de $\Aut(K)$, que recibe el nombre de \underline{grupo de automorfismos de $F\leq K$}.\\
\end{definicion}

\begin{ejercicio}\label{ej:aut_primo}
    Si $\Pi$ es el subcuerpo primo de $K$, entonces $\Aut_\Pi(K) = \Aut(K)$.\\

    \noindent
    Basta probar la inclusión $\supseteq)$. Para ello, sea $\sigma\in \Aut(K)$, si tomamos $a\in \Pi$ y $b\in K$, tenemos que:
    \begin{equation*}
        \sigma(a\cdot b) = \sigma(a)\cdot \sigma(b) \AstIg \iota(a)\cdot \sigma(b) = a\cdot \sigma(b)
    \end{equation*}
    Donde en $(\ast)$ hemos usado que $\sigma\big|_\Pi=\iota$, donde $\iota:\Pi\to K$ es la aplicación inclusión, algo que probamos en el Ejercicio~\ref{ej:homomorfismo_primo}, con lo que $\sigma\in \Aut_\Pi(K)$.
\end{ejercicio}

\begin{prop}\label{prop:comienzo_chap_2}
    Si $F\leq K$ es finita, entonces $|\Aut_F(K)| \leq [K:F]$
    \begin{proof}
        Si llamamos $F\stackrel{\iota}{\to}K$ al homomorfismo inclusión, entonces:
        \begin{equation*}
            \Aut_F(K) = \Ex(\iota, \iota)
        \end{equation*}
        \begin{description}
            \item [$\subseteq )$] Si $\sigma\in \Aut_F(K)$, tenemos entonces que $\sigma$ es $F-$lineal, y por el Ejercicio~\ref{ej:F-lineal}, tenemos entonces que $\sigma\circ\iota = \iota$, lo que nos dice que $\sigma\in \Ex(\iota,\iota)$.
            \item [$\supseteq )$] Si tomamos $\sigma\in \Ex(\iota,\iota)$ como es homomorfismo de cuerpos tenemos que es inyectivo, y la condición $\sigma\iota = \iota$ nos dice por el Ejercicio~\ref{ej:F-lineal} que $\sigma$ es $F-$lineal. Como está definida entre dos espacios vectoriales de dimensión finita, ha de ser necesariamente sobreyectivo, con lo que $\sigma\in \Aut_F(K)$
        \end{description}
        Finalmente, la segunda proposición de extensión nos dice que:
        \begin{equation*}
            |\Aut_F(K)| = |\Ex(\iota,\iota)| \leq [K:F]
        \end{equation*}
    \end{proof}
\end{prop}

\begin{notacion}
    En una situación como la de la Proposición anterior, es decir, siempre que tengamos $F\leq K$ con $\iota:F\to K$ el homomorfismo inclusión, llamaremos al conjunto $\Ex(\iota,\iota)$ extensiones de la inclusión.
\end{notacion}

\begin{prop}\label{prop:cuerpo_desc_aut_f_lin}
    Si $F\leq K$ es cuerpo de descomposición de $f\in F[x]$, entonces: 
    \begin{equation*}
        |\Aut_F(K)| \leq [K:F]
    \end{equation*}
    y si todas las raíces de $f$ en $K$ son simples (es decir, $f$ tiene $\deg f$ raíces distintas), entonces:
    \begin{equation*}
        |\Aut_F(K)| = [K:F]
    \end{equation*}
    \begin{proof}
        Si $F\leq K$ es un cuerpo de descomposición de $f\in F[x]$, tenemos entonces que si $\alpha_1, \ldots, \alpha_s$ son las raíces de $f$ en $K$ entonces $K = F(\alpha_1, \ldots, \alpha_s)$ es una extensión algebraica y finitamente generada, luego finita, de donde aplicando la Proposición anterior tenemos que $|\Aut_F(K)|\leq [K:F]$.\\

        \noindent
        Si ahora tenemos que todas las raíces de $f$ en $K$ son simples, aplicando la igualdad de la demostración anterior $|\Aut_F(K)| = |\Ex(\iota,\iota)|$ para $\iota:F\hookrightarrow K$ la aplicación inclusión, tenemos por la tercera Proposición de extensión que:
        \begin{equation*}
            |\Aut_F(K)| = |\Ex(\iota,\iota)| = [K:F]
        \end{equation*}
    \end{proof}
\end{prop}

\begin{ejemplo}\label{ejemplo:aut_sqrt3_2}
    Por el Ejercicio~\ref{ej:aut_primo}, como $\mathbb{Q}$ es el subcuerpo primo de $\mathbb{Q}\left(\sqrt[3]{2}, w\right)$ con $w$ una raíz cúbica primitiva de la unidad, tenemos que:
    \begin{equation*}
        \Aut\left(\mathbb{Q}\left(\sqrt[3]{2}, w\right)\right) = \Aut_\mathbb{Q}\left(\mathbb{Q}\left(\sqrt[3]{2}, w\right)\right)
    \end{equation*}
    
    Puesto que $\mathbb{Q}\left(\sqrt[3]{2}, w\right)$ es cuerpo de descomposición de $x^3-2\in \mathbb{Q}[x]$ y todas sus raíces son distintas, la Proposición~\ref{prop:cuerpo_desc_aut_f_lin} nos dice que:
    \begin{equation*}
        \left|\Aut_\mathbb{Q}\left(\mathbb{Q}\left(\sqrt[3]{2},w\right)\right)\right| = [\mathbb{Q}\left(\sqrt[3]{2},w\right):\mathbb{Q}] = 6
    \end{equation*}
    Por Álgebra II, tenemos que este grupo es isomorfo a $C_6$ o a $S_3$, pero por el comentario final del ejemplo de la página~\pageref{ejemplo:extension_homomorfismos}, vimos que:
    \begin{equation*}
        \Aut\left(\mathbb{Q}\left(\sqrt[3]{2},w\right)\right) = \{\eta_{j,k} : j \in \{0,1,2\}, k\in \{1,2\}\}
    \end{equation*}
    donde:
    \begin{equation*}
        \left\{\begin{array}{ccl}
                \eta_{j,k}\left(\sqrt[3]{2}\right) &=& w^j\sqrt[3]{2} \\
                \eta_{j,k}(w) &=& w^k
        \end{array}\right.
    \end{equation*}
    
    Resulta por tanto que tenemos un grupo no conmutativo:
    \begin{align*}
        \sqrt[3]{2} &\stackrel{\eta_{1,1}}{\longmapsto} w\sqrt[3]{2} \stackrel{\eta_{1,0}}{\longmapsto} w\sqrt[3]{2} \\
        \sqrt[3]{2} &\stackrel{\eta_{1,0}}{\longmapsto} w\sqrt[3]{2} \stackrel{\eta_{1,1}}{\longmapsto} w^2\sqrt[3]{2} 
    \end{align*}
    Por tanto, $\Aut\left(\mathbb{Q}\left(\sqrt[3]{2},w\right)\right)\cong S_3$.
    % // TODO: Luego sabremos que nunca podrá ser C6
\end{ejemplo}


\begin{ejercicio}\label{ej:subgrupo_finito}
    Demostrar que, si $F$ es un cuerpo, entonces cualquier subgrupo finito de $F^\times$ es cíclico. Deducimos que, en particular, $\bb{F}_q^\times$ es un grupo cíclico de orden $q-1$. (Pista: usar la descomposición cíclica de un grupo finito abeliano).\\

    \noindent
    Sea $G$ un subgrupo finito de $F^\times$, tomamos la descomposición cíclica de $G$:
    \begin{equation*}
        G \cong C_1 \oplus \ldots \oplus C_t
    \end{equation*}
    Con $C_i$ cíclico para cada $i \in \{1,\ldots,t\}$, con $|C_{i+1}|\mid |C_i|$. Sea $m = |C_1|$, para todo $g\in G$ consideramos su descomposición $g \equiv (g_1,\ldots,g_t)$ con $g_i\in C_i$ para cada $i \in \{1,\ldots,t\}$. Como el orden de $g_i$ divide a $|C_i|$, y $|C_i|\mid m$, el orden de $g_i$ divide a $m$, con lo que $g_i^m = 1$ para cada $i \in \{1,\ldots,t\}$, de donde:
    \begin{equation*}
        g^m = 1\qquad \forall g\in G
    \end{equation*}

    De esta forma, cada elemento de $G$ es raíz de $x^m-1\in F[x]$, que a lo mucho tiene $m$ raíces, con lo que $|G| \leq m\leq |G|$, de donde $|G| = m$, por lo que todos los grupos cíclicos en los que $G$ se descompone son triviales salvo $C_1$, de donde $G$ es cíclico.
\end{ejercicio}
\begin{observacion}
    Para $\bb{F}_q$, $\bb{F}_q^\times$ es un grupo cíclico de orden $q-1$.\newline
    A cualquier generador $a$ de $\bb{F}_q^\times$ se le llama elemento primitivo de $\bb{F}_q$, por lo que:
    \begin{equation*}
        \bb{F}_q = \{0, 1, a, \ldots, a^{q-2}\}
    \end{equation*}
    Por lo que $\bb{F}_q = \bb{F}_p(a)$, con $p = \car(\bb{F}_q)$.
\end{observacion}


\begin{teo}
    Sea $\bb{F}_q$ un cuerpo finito con $q=p^n$ elementos, entonces $\Aut(\bb{F}_q)$ es un grupo cíclico de orden $n$.
    \begin{proof}
        Sabemos por la Proposición~\ref{prop:cuerpo_descomposicion} que $\bb{F}_q$ es cuerpo de descomposición de $x^q-x\in \bb{F}_q[x]$, así como que las raíces de dicho polinomio son todas distintas (puesto que no comparte raíces con su polinomio derivado). Estamos en las condiciones de aplicar la Proposición~\ref{prop:cuerpo_desc_aut_f_lin}, obteniendo que:
        \begin{equation*}
            |\Aut(\bb{F}_q)| = |\Aut_{\bb{F}_p}(\bb{F}_q)| = [\bb{F}_q : \bb{F}_p]
        \end{equation*}

        Como dicho valor es $\dim_{\bb{F}_p}\bb{F}_q$, tenemos que $|\Aut(\bb{F}_q)| = n$. Sea ahora:
        \Func{\tau}{\bb{F}_q}{\bb{F}_q}{a}{a^{p}}
        
        Por el Ejercicio~\ref{ej:cuerpo_caracteristica}, se trata de un homomorfismo de cuerpos, luego $\tau \in \Aut(\bb{F}_q)$. Veamos que su orden es $n$. Para ello, sea $m\in \mathbb{N}\setminus \{0\}$ de forma que:
        \begin{equation*}
            \tau^m = \id_{\bb{F}_q}
        \end{equation*}
        En el Ejercicio~\ref{ej:subgrupo_finito} vimos que $\bb{F}_q^{\times}$ es cíclico y de orden $q-1$. Tomamos $a$ como su generador, que será de orden $q-1$, lo que nos dice entonces que:
        \begin{equation*}
            a = \id_{\bb{F}_q}(a) = \tau^m (a)  = a^{p^m}\Longrightarrow a^{p^m-1} = 1
        \end{equation*}
        Como $O(a)=q-1$, tenemos que $(p^n-1)\mid (p^m-1)$, de donde $p^n-1\leq p^m-1$, y por tanto $n\leq m$. Por otro lado, tenemos que:
        \begin{equation*}
            \tau^n(x) = x^{p^n} = x\cdot x^{p^n-1} = x\cdot 1 = x \qquad \forall x\in \bb{F}_q
        \end{equation*}
        donde se ha hecho uso de que $O(x)\mid (q-1)$, con lo que $\tau^n = \id_{\bb{F}_q}$. Como hemos visto que $m\geq n$ y que $\tau^n = \id_{\bb{F}_q}$, tenemos que el orden de $\tau$ es $n$. Por tanto, $\Aut(\bb{F}_q)$ ha de ser cíclico.
    \end{proof}
\end{teo}

\begin{definicion}[Automorfismo de Frobenius]\label{def:automorfismo_frobenius}
    Sea $\bb{F}_q$ un cuerpo finito con $q=p^n$ elementos. El automorfismo $\tau\in \Aut(\bb{F}_q)$ definido por:
    \begin{equation*}
        \tau(a) = a^p \qquad \forall a\in \bb{F}_q
    \end{equation*}
    recibe el nombre de \underline{automorfismo de Frobenius}.
\end{definicion}

Por la demostración del teorema anterior, tenemos que $\Aut(\bb{F}_q) = \langle \tau\rangle$, con $\tau$ el automorfismo de Frobenius.

\begin{comment}
    \section{Clausura Algebraica}

% // TODO: Redactado por Claude
% // De sus apuntes, solo ha dado la PROPOSICION´ 2.26.

En las secciones anteriores, dado un polinomio $f\in F[x]$, hemos construido un cuerpo donde $f$ tiene todas sus raíces (un cuerpo de descomposición). Ahora damos un paso más: buscamos un único cuerpo que sirva para \emph{todos} los polinomios de $F[x]$ a la vez. A ese cuerpo lo llamaremos clausura algebraica de $F$.


\begin{definicion}
    Sea $\sigma:F\to K$ un homomorfismo de cuerpos.
    \begin{itemize}
        \item Diremos que $\sigma$ es \underline{finito} si $K$ tiene dimensión finita como $F-$espacio vectorial (a través de $\sigma$), es decir, si la extensión $\sigma(F)\leq K$ es finita.
        \item Diremos que $\sigma$ es \underline{algebraico} si la extensión $\sigma(F)\leq K$ es algebraica.
    \end{itemize}
\end{definicion}

\begin{observacion}
    Es solo una cuestión de lenguaje: si $F\leq K$ es una extensión y $\iota:F\hookrightarrow K$ es la inclusión, que $\iota$ sea finito (o algebraico) significa exactamente que la extensión $F\leq K$ es finita (o algebraica). Recordemos además que toda extensión finita es algebraica (una extensión es finita si, y solo si, es algebraica y finitamente generada), de modo que \emph{todo homomorfismo finito es algebraico}. El recíproco es falso: $\mathbb{Q}\leq\overline{\mathbb{Q}}$ es algebraica pero no finita.
\end{observacion}

\begin{lema}[Transitividad de las extensiones algebraicas]\label{lema:transitividad_alg}
    Sean $F\leq K\leq L$ extensiones de cuerpos. Si $F\leq K$ y $K\leq L$ son algebraicas, entonces $F\leq L$ es algebraica.
    \begin{proof}
        Sea $\alpha\in L$. Como $K\leq L$ es algebraica, existe un polinomio no nulo
        \begin{equation*}
            g = \sum_{i=0}^{n}a_ix^i\in K[x], \qquad g(\alpha)=0
        \end{equation*}
        Los coeficientes $a_0,\ldots,a_n\in K$ son algebraicos sobre $F$, por serlo $F\leq K$. Por tanto, $F\leq F(a_0,\ldots,a_n)$ es finita, al estar generada por elementos algebraicos. Además, $\alpha$ es algebraico sobre $F(a_0,\ldots,a_n)$, porque es raíz de $g\in F(a_0,\ldots,a_n)[x]$, con lo que $F(a_0,\ldots,a_n)\leq F(a_0,\ldots,a_n)(\alpha)$ es finita. Por el Lema de la Torre, $F\leq F(a_0,\ldots,a_n,\alpha)$ es finita, luego algebraica, y en particular $\alpha$ es algebraico sobre $F$.
    \end{proof}
\end{lema}

\subsection*{Primer problema: juntar dos extensiones finitas}

Empezamos con una pregunta modesta. Si tenemos dos extensiones finitas $F\leq L_1$ y $F\leq L_2$, ¿existe un cuerpo $E$ que contenga ``copias'' de $L_1$ y de $L_2$ que \emph{coincidan sobre $F$}? La respuesta es afirmativa.

\begin{prop}\label{prop:fusion_finitas}
    Dados homomorfismos finitos $\sigma_1:F\to L_1$ y $\sigma_2:F\to L_2$, existen homomorfismos finitos $\tau,\tau_1,\tau_2$ que hacen conmutar el diagrama:
    \begin{figure}[H]
        \centering
        \shorthandoff{""}
        \begin{tikzcd}
        F \arrow[r, "\sigma_1"] \arrow[d, "\sigma_2"'] \arrow[rd, "\tau" description] & L_1 \arrow[d, "\tau_1"] \\
        L_2 \arrow[r, "\tau_2"']                                                      & E
        \end{tikzcd}
        \shorthandon{""}
    \end{figure}
    \noindent
    es decir, $\tau_1\sigma_1 = \tau = \tau_2\sigma_2$. Además, $\tau:F\to E$ puede tomarse como un cuerpo de descomposición de un polinomio de $F[x]$.
    \begin{proof}
        \underline{Paso 1: un polinomio para cada $L_i$.} Fijemos $i\in\{1,2\}$. Como $\sigma_i(F)\leq L_i$ es finita, tenemos que $L_i=\sigma_i(F)(\alpha_1,\ldots,\alpha_t)$ con $\alpha_1,\ldots,\alpha_t$ algebraicos sobre $\sigma_i(F)$ (basta tomar una base). Sean $g_j=\Irr(\alpha_j,\sigma_i(F))\in\sigma_i(F)[x]$. Como $\sigma_i:F\to\sigma_i(F)$ es un isomorfismo, existe $f_i\in F[x]$ con $f_i^{\sigma_i}=g_1\cdots g_t$.

        Por el último ejercicio de la Sección 1 de este tema (tomar un cuerpo de descomposición de $f_i^{\sigma_i}$ y extender homomorfismos una raíz tras otra), existe un homomorfismo de cuerpos $\lambda_i:L_i\to K_i$ tal que
        \begin{equation*}
            \lambda_i\sigma_i:F\to K_i \text{\ es un cuerpo de descomposición de\ } f_i
        \end{equation*}

        \underline{Paso 2: un cuerpo común.} Por la Proposición~\ref{prop:existe_cdd}, existe un cuerpo de descomposición $\tau:F\to E$ del polinomio $f_1f_2$. En particular, $\tau$ es finito (es algebraico y finitamente generado).

        \underline{Paso 3: encajar cada $K_i$ en $E$.} Como $f_i\mid f_1f_2$ y $(f_1f_2)^\tau$ es producto de factores lineales en $E[x]$, también $f_i^\tau$ lo es (los factores de un producto de lineales son producto de algunos de ellos, por ser $E[x]$ un DFU). Como $\lambda_i\sigma_i$ es un cuerpo de descomposición de $f_i$, la Proposición~\ref{prop:extension_3} (con $\tau$ y $\lambda_i\sigma_i$) nos dice que $\Ex(\tau,\lambda_i\sigma_i)\neq\emptyset$. Existe, por tanto, $\eta_i:K_i\to E$ con $\eta_i\lambda_i\sigma_i=\tau$.

        \underline{Paso 4: conclusión.} Tomamos $\tau_i=\eta_i\lambda_i:L_i\to E$. Entonces $\tau_i\sigma_i=\eta_i\lambda_i\sigma_i=\tau$ para $i=1,2$. Falta ver que $\tau_i$ es finito: tenemos $\tau(F)\leq\tau_i(L_i)\leq E$, y como $\tau(F)\leq E$ es finita, también lo es $\tau_i(L_i)\leq E$ (un subespacio de dimensión finita... o más precisamente, $[E:\tau_i(L_i)]$ divide a $[E:\tau(F)]$ por el Lema de la Torre).
    \end{proof}
\end{prop}

\begin{observacion}
    Interpretemos la proposición. El cuerpo $E$ contiene a $\tau(F)\cong F$ y, dentro de él, los subcuerpos $\tau_1(L_1)\cong L_1$ y $\tau_2(L_2)\cong L_2$, ambos conteniendo a $\tau(F)$. Es decir, hemos conseguido ``meter'' a la vez $L_1$ y $L_2$ en un mismo cuerpo, compatiblemente con $F$.

    El resultado se extiende sin dificultad a una cantidad \emph{finita} de homomorfismos finitos $\sigma_i:F\to L_i$, $i=1,\ldots,n$ (basta tomar un cuerpo de descomposición de $f_1\cdots f_n$). Para muchos propósitos esto es suficiente. Sin embargo, es natural preguntarse si existe una extensión algebraica de $F$ que contenga a \emph{todas} sus extensiones finitas a la vez. La respuesta es que sí, pero para probarlo hace falta inducción transfinita (el Lema de Zorn). La clave es la noción de clausura algebraica.
\end{observacion}

\subsection*{Clausura algebraica y cuerpos algebraicamente cerrados}

\begin{definicion}[Clausura algebraica]
    Sea $F$ un cuerpo. Diremos que un homomorfismo de cuerpos algebraico $\sigma:F\to C$ es una \underline{clausura algebraica} de $F$ si todo homomorfismo finito $\tau:C\to E$ es un isomorfismo.
\end{definicion}

La idea es que $C$ es una extensión algebraica de $F$ que ya no se puede agrandar mediante extensiones finitas: cualquier extensión finita de $C$ es, en realidad, el propio $C$. Intuitivamente, $C$ ya contiene todas las raíces de todos los polinomios. El siguiente resultado lo formaliza.

\begin{prop}\label{prop:clausura_equiv}
    Las siguientes condiciones son equivalentes para un homomorfismo algebraico de cuerpos $\sigma:F\to C$:
    \begin{enumerate}
        \item $\sigma:F\to C$ es una clausura algebraica de $F$.
        \item Todo polinomio no constante de $C[x]$ tiene una raíz en $C$.
    \end{enumerate}
    \begin{proof}
        \begin{description}
            \item [$1\Longrightarrow 2)$] Sea $f\in C[x]$ no constante. Tomamos un cuerpo de descomposición $\tau:C\to K$ de $f$, que es un homomorfismo finito. Como $\sigma$ es una clausura algebraica, $\tau$ es un isomorfismo. Sea $\alpha\in K$ una raíz de $f^\tau$ (existe, pues $f^\tau$ es producto de factores lineales y no constante). Aplicando $\tau^{-1}$ a $f^\tau(\alpha)=0$ obtenemos $f(\tau^{-1}(\alpha))=0$, es decir, $\tau^{-1}(\alpha)\in C$ es raíz de $f$.
            \item [$2\Longrightarrow 1)$] Supongamos que todo polinomio no constante de $C[x]$ tiene una raíz en $C$ y sea $\tau:C\to E$ un homomorfismo finito. Hemos de ver que $\tau$ es sobreyectivo (como es un homomorfismo de cuerpos, será entonces biyectivo).

            Sea $x\in E$. Como $\tau(C)\leq E$ es finita, $x$ es algebraico sobre $\tau(C)$, y como $\tau:C\to\tau(C)$ es un isomorfismo, existe un polinomio mónico $f\in C[x]$ con $f^\tau=\Irr(x,\tau(C))$. Aplicando repetidamente la hipótesis (extraemos una raíz $c_1$, dividimos por $x-c_1$, y repetimos con el cociente), $f$ se factoriza en $C[x]$ como:
            \begin{equation*}
                f=\prod_{i=1}^{n}(x-c_i), \qquad c_1,\ldots,c_n\in C
            \end{equation*}
            Por tanto, $f^\tau=\prod_{i=1}^{n}(x-\tau(c_i))$. Como $x$ es raíz de $f^\tau$ y $E$ es un cuerpo, ha de ser $x=\tau(c_i)$ para algún $i$, luego $x\in\tau(C)$. Así, $\tau$ es sobreyectivo.\qedhere
        \end{description}
    \end{proof}
\end{prop}

\begin{definicion}[Cuerpo algebraicamente cerrado]
    Un cuerpo $C$ se dice \underline{algebraicamente cerrado} si todo polinomio no constante de $C[x]$ tiene una raíz en $C$. Equivalentemente (por el argumento de la demostración anterior), si todo polinomio no constante de $C[x]$ se descompone como producto de polinomios de grado 1 en $C[x]$.
\end{definicion}

\noindent
Con esta definición, la Proposición~\ref{prop:clausura_equiv} se reformula así: \emph{una clausura algebraica de $F$ es un homomorfismo algebraico $\sigma:F\to C$ tal que $C$ es algebraicamente cerrado.}

\begin{ejemplo}
    \begin{itemize}
        \item $\mathbb{R}$ \emph{no} es algebraicamente cerrado, pues $x^2+1$ no tiene raíces reales.
        \item El Teorema Fundamental del Álgebra, que demostraremos en el tema de extensiones de Galois, afirma que $\mathbb{C}$ es algebraicamente cerrado. Como $\mathbb{R}\leq\mathbb{C}$ es finita (de grado 2), luego algebraica, la inclusión $\mathbb{R}\hookrightarrow\mathbb{C}$ es una clausura algebraica de $\mathbb{R}$.
        \item $\mathbb{Q}\leq\mathbb{C}$ \emph{no} es una clausura algebraica de $\mathbb{Q}$ porque no es algebraica ($\pi$ es trascendente). Sin embargo, sí lo es $\mathbb{Q}\leq\overline{\mathbb{Q}}$, donde $\overline{\mathbb{Q}}$ es el cuerpo de los números algebraicos. En efecto, la extensión es algebraica por definición. Veamos que $\overline{\mathbb{Q}}$ es algebraicamente cerrado. Sea $f\in\overline{\mathbb{Q}}[x]$ no constante. Como $\mathbb{C}$ es algebraicamente cerrado, existe $z\in\mathbb{C}$ con $f(z)=0$. Entonces $z$ es algebraico sobre $\overline{\mathbb{Q}}$, y como $\mathbb{Q}\leq\overline{\mathbb{Q}}$ es algebraica, el Lema~\ref{lema:transitividad_alg} nos dice que $z$ es algebraico sobre $\mathbb{Q}$, es decir, $z\in\overline{\mathbb{Q}}$.
        \item Ningún cuerpo finito $\bb{F}_q$ es algebraicamente cerrado. En efecto, el polinomio
            \begin{equation*}
                f=\left(\prod_{a\in\bb{F}_q}(x-a)\right)+1\in\bb{F}_q[x]
            \end{equation*}
            es no constante y $f(a)=1\neq 0$ para todo $a\in\bb{F}_q$. En particular, la clausura algebraica de $\bb{F}_q$ es siempre un cuerpo infinito.
    \end{itemize}
\end{ejemplo}

\subsection*{Existencia de la clausura algebraica}

Nuestro objetivo es demostrar que todo cuerpo tiene una clausura algebraica. La demostración usa el Lema de Zorn y necesita unas precisiones conjuntistas previas.

\begin{observacion}[Lema de Zorn]
    Sea $(\Gamma,\preceq)$ un conjunto parcialmente ordenado y no vacío tal que toda cadena (subconjunto totalmente ordenado) de $\Gamma$ tiene una cota superior en $\Gamma$. Entonces $\Gamma$ tiene al menos un elemento maximal, es decir, un $m\in\Gamma$ tal que si $m\preceq x$ entonces $x=m$.
\end{observacion}

\begin{definicion}[Trasladar la estructura de cuerpo]
    Sea $K$ un cuerpo y sea $\chi:K\to X$ una biyección con un conjunto cualquiera $X$. Definimos en $X$ las operaciones:
    \begin{equation*}
        x+y=\chi\left(\chi^{-1}(x)+\chi^{-1}(y)\right), \qquad x\cdot y=\chi\left(\chi^{-1}(x)\cdot\chi^{-1}(y)\right) \qquad \forall x,y\in X
    \end{equation*}
    con $0_X=\chi(0)$ y $1_X=\chi(1)$. Como $\chi$ es biyectiva, todos los axiomas de cuerpo se verifican en $X$ (se ``copian'' desde $K$) y, por construcción, $\chi:K\to X$ es un isomorfismo de cuerpos. Diremos que hemos \underline{trasladado a $X$ la estructura de cuerpo de $K$} mediante $\chi$.
\end{definicion}

\begin{observacion}
    En las notas de la asignatura esto se formula con la ``aplicación estructura'' $e:K\times K\to K\times K$, $e(x,y)=(x+y,xy)$, de modo que $\sigma:K\to L$ es homomorfismo de cuerpos si, y solo si, $\sigma(1)=1$ y $(\sigma\times\sigma)\circ e=f\circ(\sigma\times\sigma)$, donde $f$ es la aplicación estructura de $L$. Trasladar la estructura mediante $\chi$ es definir $s=(\chi\times\chi)\circ e\circ(\chi^{-1}\times\chi^{-1})$. Es lo mismo que acabamos de escribir con las operaciones explícitas.
\end{observacion}

El último ingrediente es controlar el cardinal de una extensión algebraica, para poder asegurar que ``cabe'' en un conjunto suficientemente grande.

\begin{lema}\label{lema:cardinal_alg}
    Sea $\sigma:F\to E$ un homomorfismo algebraico. Si $F$ es infinito, entonces $\#E=\#F$. Si $F$ es finito, entonces $E$ es a lo sumo numerable.
    \begin{proof}
        Comparamos primero los cardinales de $F$ y de $F[x]$. Podemos escribir $F[x]=\bigcup_{n\in\mathbb{N}}P_n$, donde $P_n$ es el $F-$espacio vectorial de los polinomios de grado menor o igual que $n$, que es isomorfo a $F^{n+1}$.
        \begin{itemize}
            \item Si $F$ es finito, cada $P_n$ es finito y $F[x]$ es una unión numerable de conjuntos finitos, luego numerable.
            \item Si $F$ es infinito, $\#P_n=\#F^{n+1}=\#F$ para todo $n$, y una unión numerable de conjuntos del mismo cardinal infinito tiene ese mismo cardinal, luego $\#F[x]=\#F$.
        \end{itemize}
        La observación clave es que, como $\sigma$ es algebraico, cada elemento de $E$ es raíz de algún $f^\sigma$ con $f\in F[x]$ no nulo. Es decir, $E$ es la unión, indexada por $F[x]\setminus\{0\}$, de los conjuntos de raíces de los $f^\sigma$ en $E$, cada uno de los cuales es finito. Por tanto, $\#E\leq\#F[x]$ (si $F[x]$ es infinito, como ocurre siempre). Si $F$ es finito, esto da que $E$ es a lo sumo numerable. Si $F$ es infinito, obtenemos $\#E\leq\#F$, y como $\sigma$ es inyectiva, $\#F\leq\#E$, de donde $\#E=\#F$.
    \end{proof}
\end{lema}

\begin{teo}[Existencia de la clausura algebraica]\label{teo:existe_clausura}
    Todo cuerpo $F$ tiene una clausura algebraica.
    \begin{proof}
        Dividimos la demostración en pasos.\\

        \noindent
        \underline{Paso 0: un conjunto ambiente $U$.} Tomamos un conjunto $U$ de cardinal no numerable y estrictamente mayor que el de $F$ (por ejemplo, el conjunto de partes de $F\cup\mathbb{N}$). Elegimos un subconjunto $F_1\subseteq U$ y una biyección $\sigma:F\to F_1$, y trasladamos a $F_1$ la estructura de cuerpo de $F$. Así, $\sigma:F\to F_1$ es un isomorfismo de cuerpos y $F_1$ es un cuerpo contenido en $U$.\\

        \noindent
        \underline{Paso 1: el conjunto $\Gamma$.} Definimos $\Gamma$ como el conjunto de los pares $(K,e)$, donde $K\subseteq U$ es un subconjunto y $e$ es una estructura de cuerpo en $K$ de modo que $F_1$ es un subcuerpo de $K$ y $F_1\leq K$ es una extensión algebraica. Ordenamos $\Gamma$ declarando:
        \begin{equation*}
            (K,e)\preceq(L,f) \Longleftrightarrow K \text{\ es un subcuerpo de\ } L
        \end{equation*}
        (es decir, $K\subseteq L$ y $f$ restringida a $K$ es $e$). Claramente es un orden parcial, y $\Gamma\neq\emptyset$ porque $(F_1,\ \text{su estructura})\in\Gamma$.\\

        \noindent
        \underline{Paso 2: toda cadena tiene cota superior.} La cadena vacía está acotada por $F_1$. Sea ahora $\cc{C}$ una cadena no vacía de $\Gamma$, y sea:
        \begin{equation*}
            M=\bigcup_{(C,c)\in\cc{C}}C\subseteq U
        \end{equation*}
        Definimos suma y producto en $M$ así: dados $x,y\in M$, como $\cc{C}$ está totalmente ordenada, existe $(C,c)\in\cc{C}$ con $x,y\in C$ (basta tomar el mayor de los dos elementos de la cadena a los que pertenecen $x$ e $y$), y definimos $x+y$ y $xy$ como en $C$. Esto no depende del $(C,c)$ elegido: si $(C',c')$ es otro elemento de la cadena que contiene a $x,y$, será comparable con $C$, y como uno es subcuerpo del otro, las operaciones coinciden. Los axiomas de cuerpo se comprueban fácilmente, ya que cada axioma involucra solo una cantidad finita de elementos (por ejemplo tres, para la asociatividad), y cualquier conjunto finito de elementos de $M$ está contenido en un mismo $C$ de la cadena. El cero y el uno son los de $F_1$. Además, $M$ es una extensión algebraica de $F_1$, porque cada elemento de $M$ pertenece a algún $C$, y $F_1\leq C$ es algebraica. Por tanto, $(M,m)\in\Gamma$ es una cota superior de $\cc{C}$.\\

        \noindent
        \underline{Paso 3: elemento maximal.} Por el Lema de Zorn, $\Gamma$ tiene un elemento maximal $(\overline{F},t)$. Vamos a ver que $\overline{F}$ es una clausura algebraica de $F$. Consideramos el isomorfismo $\sigma:F\to F_1$ seguido de la inclusión $F_1\subseteq\overline{F}$, que nos da un homomorfismo de cuerpos, que seguimos llamando $\sigma:F\to\overline{F}$. Como $(\overline{F},t)\in\Gamma$, la extensión $\sigma(F)=F_1\leq\overline{F}$ es algebraica, es decir, $\sigma$ es algebraico.\\

        \noindent
        \underline{Paso 4: $\overline{F}$ no se puede agrandar con extensiones finitas.} Sea $\tau:\overline{F}\to E$ un homomorfismo finito arbitrario. Hemos de probar que es un isomorfismo, es decir, que es sobreyectivo.

        Por el Lema~\ref{lema:cardinal_alg}, el cardinal de $\overline{F}$ es el de $F$ (si $F$ es infinito) o es a lo sumo numerable (si $F$ es finito). Como $\tau$ es finito, luego algebraico, el mismo lema aplicado a $\tau$ nos da que $E$ tiene el mismo cardinal que $\overline{F}$ o es a lo sumo numerable. En ambos casos, $\#\overline{F}$ y $\#E$ son estrictamente menores que $\#U$. Por tanto, el conjunto $U\setminus\overline{F}$ tiene cardinal $\#U$ y existe una aplicación biyectiva
        \begin{equation*}
            b:E\setminus\tau(\overline{F})\longrightarrow X
        \end{equation*}
        para cierto subconjunto $X\subseteq U$ disjunto de $\overline{F}$.

        Definimos ahora la biyección:
        \begin{equation*}
            \chi:E\longrightarrow\overline{F}\cup X, \qquad \chi(y)=\left\{\begin{array}{ll}
                \tau^{-1}(y) & \text{si\ } y\in\tau(\overline{F}) \\
                b(y) & \text{si\ } y\notin\tau(\overline{F})
            \end{array}\right.
        \end{equation*}
        y trasladamos a $\overline{F}\cup X$ la estructura de cuerpo de $E$ mediante $\chi$, obteniendo una estructura $s$. Entonces $\chi:E\to\overline{F}\cup X$ es un isomorfismo de cuerpos y la composición
        \begin{equation*}
            \chi\tau:\overline{F}\longrightarrow\overline{F}\cup X
        \end{equation*}
        es un homomorfismo de cuerpos (composición de homomorfismos). Pero, por la definición de $\chi$, $\chi\tau$ es precisamente la \emph{inclusión} $\overline{F}\subseteq\overline{F}\cup X$. Que la inclusión sea homomorfismo significa justamente que $\overline{F}$ es un subcuerpo de $\overline{F}\cup X$ con la estructura $s$, es decir, que $s$ extiende a $t$.

        Además, $\overline{F}\leq\overline{F}\cup X$ es finita: su estructura de $\overline{F}-$espacio vectorial se obtiene trasladando, mediante el isomorfismo $\chi$, la de $\tau(\overline{F})-$espacio vectorial de $E$, que tiene dimensión finita. Como $F_1\leq\overline{F}$ es algebraica y $\overline{F}\leq\overline{F}\cup X$ es finita (luego algebraica), el Lema~\ref{lema:transitividad_alg} nos dice que $F_1\leq\overline{F}\cup X$ es algebraica. En definitiva, $(\overline{F}\cup X,s)\in\Gamma$ y $(\overline{F},t)\preceq(\overline{F}\cup X,s)$.

        Como $(\overline{F},t)$ es maximal, ha de ser $\overline{F}\cup X=\overline{F}$, es decir, $X=\emptyset$. Como $b$ es una biyección sobre $X$, esto implica que $E\setminus\tau(\overline{F})=\emptyset$, luego $\tau$ es sobreyectivo y, por tanto, un isomorfismo.\qedhere
    \end{proof}
\end{teo}

\subsection*{Unicidad de la clausura algebraica}

Resta estudiar en qué sentido es única. Primero demostramos que una clausura algebraica es ``universal'': toda extensión algebraica de $F$ se sumerge en ella.

\begin{prop}\label{prop:ext_a_clausura}
    Sea $\tau:F\to E$ un homomorfismo algebraico y sea $\sigma:F\to C$ una clausura algebraica de $F$. Entonces existe un homomorfismo de cuerpos $\eta:E\to C$ tal que $\eta\tau=\sigma$.
    \begin{figure}[H]
        \centering
        \shorthandoff{""}
        \begin{tikzcd}
        F \arrow[r, "\sigma"] \arrow[d, "\tau"'] & C \\
        E \arrow[ru, "\eta"', dashed]            &  
        \end{tikzcd}
        \shorthandon{""}
    \end{figure}
    \begin{proof}
        Aplicamos el Lema de Zorn al conjunto $\Lambda$ de todos los pares $(L,\lambda)$, donde $L$ es un cuerpo intermedio $\tau(F)\leq L\leq E$ y $\lambda:L\to C$ es un homomorfismo de cuerpos tal que $\lambda\tau=\sigma$ (aquí $\tau$ se ve como $\tau:F\to L$). El orden es:
        \begin{equation*}
            (L,\lambda)\preceq(K,\kappa)\Longleftrightarrow L\subseteq K \text{\ y\ } \kappa\big|_L=\lambda
        \end{equation*}
        El conjunto $\Lambda$ es no vacío, pues contiene a $(\tau(F),\sigma\tau^{-1})$, donde $\tau^{-1}:\tau(F)\to F$ es el inverso del isomorfismo $\tau:F\to\tau(F)$. Además, toda cadena de $\Lambda$ tiene cota superior: se toma la unión de los cuerpos $L$ de la cadena, y se define $\lambda$ en un elemento como el valor que le asigna cualquiera de los $\lambda_L$ de la cadena que lo contenga (esto está bien definido, porque los $\lambda_L$ son compatibles entre sí). Por el Lema de Zorn, $\Lambda$ tiene un elemento maximal $(M,\eta)$.

        Veamos que $M=E$. Sea $a\in E$. Como $\tau$ es algebraico, $a$ es algebraico sobre $\tau(F)\leq M$, luego sobre $M$, y podemos considerar $f=\Irr(a,M)\in M[x]$. Entonces $f^\eta\in C[x]$ es no constante, y como $C$ es algebraicamente cerrado (Proposición~\ref{prop:clausura_equiv}), tiene una raíz $c\in C$. Aplicamos la Proposición de extensión a:
        \begin{figure}[H]
            \centering
            \shorthandoff{""}
            \begin{tikzcd}
            M \arrow[r, "\eta"] \arrow[rd, "\iota"'] & C                 \\
                                                     & M(a) \arrow[u, "\mu"', dashed]
            \end{tikzcd}
            \shorthandon{""}
        \end{figure}
        \noindent
        con $\iota$ la inclusión, el polinomio irreducible $f\in M[x]$ y $a\in M(a)$ con $f^\iota(a)=0$. Como el conjunto de raíces de $f^\eta$ en $C$ es no vacío, $\Ex(\eta,\iota)\neq\emptyset$, luego existe $\mu:M(a)\to C$ homomorfismo con $\mu\iota=\eta$, es decir, $\mu\big|_M=\eta$. En particular, $\mu\tau=\eta\tau=\sigma$, de modo que $(M(a),\mu)\in\Lambda$ y $(M,\eta)\preceq(M(a),\mu)$. Por maximalidad, $M(a)=M$, es decir, $a\in M$. Como $a\in E$ era arbitrario, $M=E$ y $\eta:E\to C$ cumple $\eta\tau=\sigma$.
    \end{proof}
\end{prop}

\begin{teo}[Unicidad de la clausura algebraica]\label{teo:unicidad_clausura}
    Si $\tau:F\to E$ y $\sigma:F\to C$ son clausuras algebraicas de un cuerpo $F$, entonces existe un isomorfismo de cuerpos $\eta:E\to C$ tal que $\eta\tau=\sigma$.
    \begin{proof}
        Por la Proposición~\ref{prop:ext_a_clausura} (como $\tau$ es algebraico y $\sigma$ es clausura algebraica), existe un homomorfismo $\eta:E\to C$ con $\eta\tau=\sigma$. Como es un homomorfismo de cuerpos, es inyectivo, y basta ver que es sobreyectivo.

        Sea $c\in C$. Como $\sigma$ es algebraico, $c$ es algebraico sobre $\sigma(F)=\eta\tau(F)\leq\eta(E)$, luego sobre $\eta(E)$. Por tanto, la extensión $\eta(E)\leq\eta(E)(c)$ es finita, y el homomorfismo
        \begin{equation*}
            E\stackrel{\eta}{\longrightarrow}\eta(E)\hookrightarrow\eta(E)(c)
        \end{equation*}
        es finito. Como $\tau:F\to E$ es una clausura algebraica, este homomorfismo ha de ser un isomorfismo; en particular, es sobreyectivo, lo que solo es posible si $\eta(E)=\eta(E)(c)$, es decir, $c\in\eta(E)$. Así, $\eta$ es sobreyectivo y, por tanto, un isomorfismo.
    \end{proof}
\end{teo}

\begin{notacion}
    Por los Teoremas~\ref{teo:existe_clausura} y~\ref{teo:unicidad_clausura}, todo cuerpo $F$ tiene una clausura algebraica, única salvo isomorfismos. Hablaremos por ello de \underline{``la'' clausura algebraica de $F$} y la notaremos por $\overline{F}$, del mismo modo que hacíamos con $\bb{F}_q$. Si $F\leq K$ y $K$ es algebraicamente cerrado, la clausura algebraica de $F$ en $K$ (el subcuerpo de los elementos de $K$ algebraicos sobre $F$) es una clausura algebraica de $F$ en el sentido de esta sección (por ejemplo, $\overline{\mathbb{Q}}\leq\mathbb{C}$).
\end{notacion}

\begin{observacion}
    La unicidad es \emph{salvo isomorfismo}, y el isomorfismo $\eta$ no tiene por qué ser único. Por ejemplo, tanto la identidad como la conjugación compleja son isomorfismos $\mathbb{C}\to\mathbb{C}$ compatibles con la inclusión de $\mathbb{R}$ (ambos fijan $\mathbb{R}$). Esto se debe a que la Proposición~\ref{prop:ext_a_clausura} no afirma unicidad, solo existencia, y de hecho $\Ex(\tau,\sigma)$ puede tener muchos elementos.
\end{observacion}
\end{comment}




% Estos ejercicios ya se han añadido antes
\begin{comment}
\begin{ejercicio} 
    Sea $F\leq K$ una extensión de cuerpos de grado 2. Mostrar que, si la característica de $F$ es distinta de dos, existe $\beta\in K$ tal que $\beta^2 \in F$ y $K = F(\beta)$.\\

    \noindent
    Sea $\alpha\in K\setminus F$, tenemos que $\alpha$ tiene grado 2 sobre $K$, puesto que si fuera de grado 1, entonces existe un polinómico mónico de grado 1 $x-a$ (con $a\in F$) de forma que $\alpha$ es raíz de dicho polinomio, con lo que ha de ser $a = \alpha\notin F$, contradicción. De esta forma, $\deg \Irr(\alpha,F)=2$, es decir, existen $a,b\in F$ de forma que $\alpha$ es raíz del polinomio:
    \begin{equation*}
        x^2+ax+b
    \end{equation*}
    Por lo que $\alpha^2 + a\cdot \alpha + b = 0$. Como la característica de $F$ no es dos, tenemos que $1+1 = 2 \neq 0$, con lo que podemos considerar $2^{-1}$. Si tomamos:
    \begin{equation*}
        \beta = \alpha + \frac{a}{2}
    \end{equation*}

    tenemos que:
    \begin{equation*}
        \beta^2 = {\left(\alpha+\frac{a}{2}\right)}^{2} = \alpha^2 + \alpha\cdot a + \frac{a^2}{4} = -b + \frac{a^{2}}{4} \in F
    \end{equation*}
    Y además $\beta\notin F$, pues $\alpha = \beta-\frac{a}{2}$. Como $\beta\in K$, es obvio que $F(\beta)\leq K$, y como $[F(\beta):F] = [K:F]$, ha de ser $K = F(\beta)$.
\end{ejercicio}

\begin{ejercicio}
    Calcular un cuerpo de descomposición de $x^4+16\in \mathbb{Q}[x]$.\\

    \noindent
    Tenemos que:
    \begin{equation*}
        x^4+16= 0 \Longleftrightarrow x=\sqrt[4]{-16} = 2\sqrt[4]{-1}
    \end{equation*}
    Si recordamos que:
    \begin{equation*}
        \sqrt[4]{-1} = \left\{e^{\frac{i}{n}(\pi + 2k\pi)} : k\in \{0,1,2,3\}\right\} = \left\{e^{\frac{i\pi}{4}}, e^{\frac{3i\pi}{4}}, e^{\frac{5i\pi}{4}}, e^{\frac{7i\pi}{4}}\right\}
    \end{equation*}

    con:
    \begin{align*}
        e^{\frac{i\pi}{4}} &= \cos\left(\frac{\pi}{4}\right) + i\sen\left(\frac{\pi}{4}\right) = \frac{\sqrt{2}}{2} + i\frac{\sqrt{2}}{2} \\
        e^\frac{3i\pi}{4} &= \cos\left(\frac{3\pi}{4}\right) + i\sen\left(\frac{3\pi}{4}\right) = -\frac{\sqrt{2}}{2} + i\frac{\sqrt{2}}{2} \\
        e^\frac{5i\pi}{4} &= \cos\left(\frac{5\pi}{4}\right) + i\sen\left(\frac{5\pi}{4}\right) = -\frac{\sqrt{2}}{2} - i\frac{\sqrt{2}}{2} \\
        e^\frac{7i\pi}{4} &= \cos\left(\frac{7\pi}{4}\right) + i\sen\left(\frac{7\pi}{4}\right) = \frac{\sqrt{2}}{2} - i\frac{\sqrt{2}}{2} 
    \end{align*}
    Por lo que:
    \begin{equation*}
        \sqrt[4]{-16} = \left\{\sqrt{2}+i\sqrt{2}, -\sqrt{2}+i\sqrt{2}, -\sqrt{2}-i\sqrt{2},\sqrt{2}-i\sqrt{2}\right\}
    \end{equation*}
    Con lo que $\mathbb{Q}\left(\sqrt{2}+i\sqrt{2}, \sqrt{2}-i\sqrt{2}\right)$ es un cuerpo de descomposición de $x^4+16$, que trataremos de probar que es igual a $\mathbb{Q}\left(i,\sqrt{2}\right)$:
    \begin{description}
        \item [$\subseteq )$] Es claro que $\mathbb{Q}\left(\sqrt{2}+i\sqrt{2}, \sqrt{2}-i\sqrt{2}\right)\leq \mathbb{Q}\left(i,\sqrt{2}\right)$.
        \item [$\supseteq )$] Vemos que:
            \begin{align*}
                \sqrt{2} &= \dfrac{\sqrt{2}+i\sqrt{2}+\sqrt{2}-i\sqrt{2}}{2} \in \mathbb{Q}\left(\sqrt{2}+i\sqrt{2}, \sqrt{2}-i\sqrt{2}\right) \\
                i &= \dfrac{\sqrt{2}+i\sqrt{2}-\sqrt{2}}{\sqrt{2}} \in  \mathbb{Q}\left(\sqrt{2}+i\sqrt{2}, \sqrt{2}-i\sqrt{2}\right)
            \end{align*}
    \end{description}
    En definitiva, $\mathbb{Q}\left(i,\sqrt{2}\right)$ es un cuerpo de descomposición de $x^4+16$.
\end{ejercicio}

\begin{ejercicio} 
    Razonar cuáles de los siguientes números complejos son algebraicos sobre $\mathbb{Q}$, suponiendo conocido que $e$ y $\pi$ son trascendentes:
    \begin{equation*}
        \sqrt[5]{4}, \left(1+\sqrt[5]{4}\right){\left(1-\sqrt[5]{16}\right)}^{-1}, \pi^2, e^2-i, i\sqrt{i}+\sqrt{2}, \sqrt{1-\sqrt[3]{2}}, \sqrt{\pi}, \sqrt{2}{\left(\sqrt[3]{2}+\sqrt[5]{2}\right)}^{-1}.
    \end{equation*}
    Veamos cada caso:
    \begin{itemize}
        \item $\sqrt[5]{4}$ es algebraico sobre $\mathbb{Q}$, puesto que es raíz de $x^5-4$.
        \item $\left(1+\sqrt[5]{4}\right){\left(1-\sqrt[5]{16}\right)}^{-1}$

            En el apartado anterior hemos visto que $[\mathbb{Q}\left(\sqrt[5]{4}\right) :\mathbb{Q}]\leq 5$, con lo que la extensión $\mathbb{Q}\leq \mathbb{Q}\left(\sqrt[5]{4}\right)$ es finita, luego algebraica y finitamente generada, por lo que todo elemento de este último cuerpo será algebraico sobre $\mathbb{Q}$. Observemos que:
            \begin{equation*}
                (1+\sqrt[5]{4}){\left(1-\sqrt[5]{16}\right)}^{-1} = (1+\sqrt[5]{4}){\left(1-\sqrt[5]{4}\sqrt[5]{4}\right)}^{-1} \in \mathbb{Q}(\sqrt[5]{4})
            \end{equation*}
            Por lo que es algebraico sobre $\mathbb{Q}$.
        \item $\pi^2$ 

            Por reducción al absurdo, si fuera $\pi^2$ algebraico sobre $\mathbb{Q}$ tendríamos entonces que existe un polinomio:
            \begin{equation*}
                f = \sum_{i = 0}^n f_i x^i, \qquad f_i \in \mathbb{Q} \quad \forall i \in \{0,\ldots,n\}
            \end{equation*}
            de forma que $f(\pi^2) = 0$. En dicho caso, si tomamos:
            \begin{equation*}
                h = \sum_{i=0}^n f_i x^{2i} \in \mathbb{Q}[x]
            \end{equation*}
            obtendríamos entonces que:
            \begin{equation*}
                h(\pi) = \sum_{i=0}^n f_i \pi^{2i} = \sum_{i=0}^n f_i {(\pi^2)}^{i} = f(\pi^2) = 0
            \end{equation*}
            y esto es una contradicción con que $\pi$ es trascendente.
        \item $e^2-i$ 

            Por reducción al absurdo, si fuera $e^2-i\in \overline{\mathbb{Q}}$, por ser $i \in \overline{\mathbb{Q}}$ (es raíz de $x^2+1$) tendríamos entonces que $e^2\in \overline{\mathbb{Q}}$, de donde seríamos capaces de razonar que $e\in \overline{\mathbb{Q}}$, al igual que en el caso $\pi^2$, llegando a una contradicción, puesto que $e$ es trascendente.
        \item $i\sqrt{i}+\sqrt{2}$ 

            Vemos que $i \sqrt{i}+\sqrt{2} \in \mathbb{Q}\left(\sqrt{i},\sqrt{2}\right)$, y aplicando el Lema de la Torre vemos que:
            \begin{equation*}
                \left[\mathbb{Q}\left(\sqrt{i},\sqrt{2}\right):\mathbb{Q}\right] = \left[\mathbb{Q}\left(\sqrt{i},\sqrt{2}\right)\right]\left[\mathbb{Q}\left(\sqrt{2}\right):\mathbb{Q}\right]
            \end{equation*}
            con:
            \begin{itemize}
                \item $\left[\mathbb{Q}\left(\sqrt{2}\right):\mathbb{Q}\right] = 2$ por ser $x^2-2=\Irr\left(\sqrt{2},\mathbb{Q}\right)$, irreducible por Eisenstein para $p=2$.
                \item $\left[\mathbb{Q}\left(\sqrt{i},\sqrt{2}\right):\mathbb{Q}\left(\sqrt{2}\right)\right]\leq 4$, ya que $x^4+1$ tiene a $\sqrt{i}$ como raíz.
            \end{itemize}
            En definitiva, la extensión $\mathbb{Q}\leq \mathbb{Q}\left(\sqrt{i},\sqrt{2}\right)$ es finita, luego algebraica, de donde el elemento $i\sqrt{i}+\sqrt{2}$ es algebraico.
        \item $\sqrt{1-\sqrt[3]{2}}$

            Buscamos un polinomio con coeficientes en $\mathbb{Q}$ del que este elemento sea raíz. Para ello, lo que haremos será ver que este ha de cumplir que:
            \begin{equation*}
                x^2 = 1-\sqrt[3]{2} \Longrightarrow x^2-1 = -\sqrt[3]{2}
            \end{equation*}
            De donde:
            \begin{equation*}
                {(x^2-1)}^{3} = x^6-3x^4+3x^2-1 = -2
            \end{equation*}
            Por lo que si tomamos $f = x^6-3x^4+3x^2+1\in \mathbb{Q}[x]$, tenemos que $\sqrt{1-\sqrt[3]{2}}$ es raíz de $f$, con lo que es algebraico sobre $\mathbb{Q}$.
        \item $\sqrt{\pi}$ 

            Si fuera $\sqrt{\pi}$ un número algebraico tendríamos entonces que la extensión $\mathbb{Q}\leq \mathbb{Q}\left(\sqrt{\pi}\right)$ sería finita, y como tenemos:
            \begin{equation*}
                \mathbb{Q}\leq \mathbb{Q}(\pi) \leq \mathbb{Q}\left(\sqrt{\pi}\right)
            \end{equation*}
            ya que $\pi = {\left(\sqrt{\pi}\right)}^{2}$, tendremos entonces por el Lema de la Torre que la extensión $\mathbb{Q}\leq \mathbb{Q}(\pi)$ es finita, luego algebraica, por lo que $\pi$ sería algebraico, contradicción que viene de suponer que $\sqrt{\pi}$ es un número algebraico.
        \item $\sqrt{2}{(\sqrt[3]{2}+\sqrt[5]{2})}^{-1}$ 

            Por el Lema de la Torre, tenemos que:
            \begin{align*}
                &\left[\mathbb{Q}\left(\sqrt{2},\sqrt[3]{2},\sqrt[5]{2}\right):\mathbb{Q}\right] =
             \\ &\left[\mathbb{Q}\left(\sqrt{2},\sqrt[3]{2},\sqrt[5]{2}\right):\mathbb{Q}\left(\sqrt{2},\sqrt[3]{2}\right)\right] \left[\mathbb{Q}\left(\sqrt{2},\sqrt[3]{2}\right):\mathbb{Q}(\sqrt{2})\right] \left[\mathbb{Q}\left(\sqrt{2}\right):\mathbb{Q}\right]
            \end{align*}
            Con:
            \begin{itemize}
                \item $\left[\mathbb{Q}\left(\sqrt{2}\right):\mathbb{Q}\right] = 2$, ya que $\Irr(\sqrt{2},\mathbb{Q}) = x^2-2$ por Eisenstein para $p=2$.
                \item $\left[\mathbb{Q}\left(\sqrt{2},\sqrt[3]{2}\right):\mathbb{Q}\left(\sqrt{2}\right)\right]\leq 3$ ya que $x^3-2$ es un polinomio con $\sqrt[3]{2}$ como raíz.
                \item $\left[\mathbb{Q}\left(\sqrt{2},\sqrt[3]{2},\sqrt[5]{2}\right):\mathbb{Q}\left(\sqrt{2},\sqrt[3]{2}\right)\right] \leq 5$, ya que $x^5-2$ es un polinomio con $\sqrt[5]{2}$ como raíz.
            \end{itemize}
            En definitiva, la extensión $\mathbb{Q}\leq \mathbb{Q}\left(\sqrt{2},\sqrt[3]{2},\sqrt[5]{2}\right)$ es finita, luego algebraica y finitamente generada. En particular, tenemos que:
            \begin{equation*}
                \sqrt{2}{\left(\sqrt[3]{2}+\sqrt[5]{2}\right)}^{-1} \in \mathbb{Q}\left(\sqrt{2},\sqrt[3]{2},\sqrt[5]{2}\right)
            \end{equation*}
            Por lo que $\sqrt{2}{\left(\sqrt[3]{2}+\sqrt[5]{2}\right)}^{-1}$ es algebraico.
    \end{itemize}
\end{ejercicio}

\begin{ejercicio}
    Sea $F\leq K$ una extensión de cuerpos, $\alpha\in K$ y $n$ natural no nulo. Demostrar que $\alpha$ es algebraico sobre $F$ si, y solo si, $\alpha^n$ es algebraico sobre $F$.\\

    \begin{description}
        \item [$\Longrightarrow )$] Si $\alpha$ es algebraico sobre $F$ entonces la extensión $F\leq F(\alpha)$ es algebraica, y tenemos que $\alpha^n \in F(\alpha)$, por lo que $\alpha^n$ es algebraico sobre $F$.
        \item [$\Longleftarrow )$] Si $\alpha^n$ es algebraico sobre $F$, entonces existe $f\in F[x]$ de forma que $f(\alpha^n) = 0$. Si $f$ se escribe como:
            \begin{equation*}
                f = \sum_{i=1}^{m}f_ix^i \qquad f_i \in F
            \end{equation*}
            tenemos entonces que:
            \begin{equation*}
                f(\alpha^n) = \sum_{i=1}^{m}f_i{(\alpha^n)}^{i} = 0
            \end{equation*}
            Por tanto, si consideramos el polinomio:
            \begin{equation*}
                g = \sum_{k=1}^{m} f_i x^{in} \in F[x]
            \end{equation*}
            tendremos entonces:
            \begin{equation*}
                g(\alpha) = \sum_{k=1}^{m} f_i \alpha^{in} = \sum_{k=1}^{m} f_i{(\alpha^n)}^{i} = 0
            \end{equation*}
            Por lo que $\alpha$ es algebraico sobre $F$.
    \end{description}
\end{ejercicio}

\begin{ejercicio}
    Sea $F\leq K$ una extensión de cuerpos, $\alpha\in K$ y $\beta = 1+\alpha^2 + \alpha^5$. Demostrar que $\alpha$ es algebraico sobre $F$ si, y solo si, $\beta$ es algebraico sobre $F$:

    \begin{description}
        \item [$\Longrightarrow )$] Si $\alpha$ es algebraico sobre $F$ entonces la extensión $F\leq F(\alpha)$ es algebraica, y tenemos que:
            \begin{equation*}
                \beta = 1+\alpha^2 + \alpha^5 \in F(\alpha)
            \end{equation*}
            Por lo que $\beta$ es algebraico sobre $F$.
        \item [$\Longleftarrow )$] Tenemos las extensiones $F\leq F(\beta)\leq F(\alpha)$. Como $\beta$ es algebraico sobre $F$ tenemos que la primera extensión es finita. Si observamos ahora que:
            \begin{equation*}
                \beta = 1 + \alpha^2 + \alpha^5 \quad\Longleftrightarrow\quad 1-\beta + \alpha^2 + \alpha^5= 0
            \end{equation*}
            Vemos que $\alpha$ es algebraico sobre $F(\beta)$, pues es raíz del polinomio $1-\beta + x^2+x^5\in F(\beta)$, por lo que la extensión $F(\beta)\leq F(\alpha)$ es también finita. El Lema de la Torre nos dice entonces que $F\leq F(\alpha)$ es finita, por lo que es algebraica, de donde $\alpha$ es algebraico sobre $F$.
    \end{description}
\end{ejercicio}
\end{comment}




\begin{ejercicio}
    Pongamos $\bb{F}_4 = \bb{F}_2(a)$ con $a^2+a+1=0$. Comprobar que $\bb{F}_{16}$ puede presentarse como $\bb{F}_{16}=\bb{F}_2(b)$, donde $b^4+b+1=0$. Determinar todos los homomorfismos de cuerpos $\bb{F}_4\to \bb{F}_{16}$ en función de $a$ y $b$.\\

    \noindent
    Veamos en primer lugar a qué nos referimos cuando decimos que $\bb{F}_4 = \bb{F}_2(a)$ con $a^2+a+1=0$. Como $x^2+x+1\in \bb{F}_2[x]$ es irreducible por no tener raíces en $\bb{F}_2$, por el Lema~\ref{lema:cuerpo_descomposicion} existe un homomorfismo $\sigma:\bb{F}_2\to K$ y $a\in K$ de forma que $K = \sigma(\bb{F}_2)(a)$ y $a^2+a+1=0$. Además, notemos que $\Irr(a,\bb{F}_2) = x^2+x+1$, por lo que $[K:\bb{F}_2] = 2$. Como $\bb{F}_2$ tiene 2 elementos, vemos entonces que $|K| = 2^2 = 4$, luego $\bb{F}_4 = K$. Por último, $\sigma(\bb{F}_2) = \bb{F}_2$, luego $\bb{F}_4 = \bb{F}_2(a)$ con $a^2+a+1=0$.

    \noindent
    Para ver que $\bb{F}_{16} = \bb{F}_2(b)$ con $b^4+b+1=0$, lo que hacemos es ver primero que el polinomio $f = x^4+x+1\in \bb{F}_2[x]$ es irreducible, ya que:
    \begin{itemize}
        \item Como $f$ no tiene raíces en $\bb{F}_2$ tenemos que no puede tener factores de grado 1, por lo que tampoco puede tenerlos de grado 3.
        \item El único factor de grado 2 que puede tener $f$ es el único polinomio irreducible de grado 2 en $\bb{F}_2[x]$, que es $x^2+x+1$, pero esto solo puede suceder si $f = {(x^2+x+1)}^{2}$, y tenemos que:
            \begin{equation*}
                {(x^2+x+1)}^{2} = x^4+x^2+1 \neq f
            \end{equation*}
            Por lo que $f$ tampoco tiene factores de grado 2, luego $f$ es irreducible.
    \end{itemize}
    Como $f$ es irreducible, existe un homomorfismo de cuerpos $\sigma:\bb{F}_2\to K$ y $b\in K$ de forma que $b^4+b+1=0$ y $K = \bb{F}_2(b)$. Como:
    \begin{equation*}
        \Irr(b,\bb{F}_2) = x^4+x+1
    \end{equation*}
    tenemos entonces que $[K:\bb{F}_2] = 4$. Como $\bb{F}_2$ tiene 2 elementos, vemos entonces que $|K| = 2^4 = 16$, luego $\bb{F}_{16} = K$.\\

    \noindent
    Para ver el número de homomorfismos de cuerpos $\eta:\bb{F}_4\to \bb{F}_{16}$ que hay, como $\bb{F}_2\leq \bb{F}_4$ y $\bb{F}_2\leq \bb{F}_{16}$, si consideramos las inclusiones en cada extensión vemos que el problema de calcular todos los homomorfismos es equivalente al problema de calcular todas las extensiones de la inclusión.

    \begin{figure}[H]
        \centering
        \shorthandoff{""}
        \begin{tikzcd}
            \mathbb{F}_2 \arrow[r, "\iota", hook] \arrow[rd, "\iota", hook] & \mathbb{F}_2(b)                    \\
                                                                            & \mathbb{F}_2(a) \arrow[u, "\eta"']
        \end{tikzcd}
        \shorthandon{""}
    \end{figure}

    \noindent
    Para ello, como $\Irr(a,\bb{F}_2) = x^2+x+1\in \bb{F}_2[x]$ y $\bb{F}_4 = \bb{F}_2(a)$, tendremos tantas extensiones de la inclusión como raíces tenga el polinomio $x^2+x+1$ en $\bb{F}_2(b)$. Buscar raíces de $x^2+x+1$ en $\bb{F}_2(b)$ es tedioso, por lo que vamos a buscar otro método. Tenemos que:
    \begin{equation*}
        x^3-1 = (x-1)(x^2+x+1)
    \end{equation*}
    Por lo que buscar raíces de $x^2+x+1$ es equivalente a buscar elementos en $\bb{F}_{16}$ distintos de 1 que cumplan $x^3 = 1$ (y por tanto cumplan que su orden en el grupo $\bb{F}_{16}^\times$ divida a $3$). Como el $3$ es primo y hemos descartado la unidad, tenemos que las raíces de $x^2+x+1$ son los elementos de orden $3$ en $\bb{F}_{16}^\times$.
    \begin{equation*}
        \{x\in \bb{F}_{16} : x^2+x+1 = 0\} = \{x\in \bb{F}_{16}^\times : O(x) = 3\}
    \end{equation*}

    Para buscar de manera cómoda los elementos de orden 3 en $\bb{F}_{16}^\times$, puesto que sabemos que $\bb{F}_{16}^\times$ es cíclico, vamos a buscar un generador de $\bb{F}_{16}^\times$. Para ello, vamos a ver que $b$ es un generador de $\bb{F}_{16}^\times$ probando que $O(b) = 15$. Como $\bb{F}_{16}^\times$ es un grupo finito, sabemos que $O(b)\mid 15$, por lo que $O(b) \in \{1,3,5,15\}$. Vamos a descartar los órdenes menores que 15.
    \begin{itemize}
        \item $O(b) = 1$ es imposible, ya que $b\neq 1$, pues $1$ no es raíz de $x^4+x+1$.
        \item $O(b) = 3$ es imposible, ya que $b^3\neq 1$, pues $\{1,b,b^2, b^3\}$ es una $\bb{F}_2-$base de $\bb{F}_{16}$, luego $b^3+1\neq 0$.
        \item $O(b) = 5$ es imposible, ya que $b^5=b^4b = (b+1)b = b^2+b \neq 1$, pues $\{1,b,b^2, b^3\}$ es una $\bb{F}_2-$base de $\bb{F}_{16}$, luego $b^5+1\neq 0$.
    \end{itemize}
    En definitiva, $O(b) = 15$, por lo que $\bb{F}_{16}^\times = \langle b \rangle $. Para buscar los elementos de orden 3, tenemos que:
    \begin{equation*}
        O(b^k) = \frac{15}{\mcd(15,k)} = 3\iff \mcd(15,k) = 5 \iff k \in \{5,10\}
    \end{equation*}

    Por tanto, los elementos de orden 3 en $\bb{F}_{16}^\times$ son $b^5$ y $b^{10}$, que son las raíces de $x^2+x+1$ en $\bb{F}_{16}$. Puesto que el enunciado nos pide dar los homomorfismos de cuerpos $\eta:\bb{F}_4\to \bb{F}_{16}$ en función de $a$ y $b$, estas raíces son:
    \begin{align*}
        b^5 &= b^4b = (b+1)b = b^2+b \\
        b^{10} &= {(b^5)}^2 = {(b^2+b)}^2 = b^4+b^2 = (b+1)+b^2 = b^2+b+1
    \end{align*}

    Por último, consideramos un homomorfismo de cuerpos $\eta:\bb{F}_4\to \bb{F}_{16}$. Como una $\bb{F}_2-$base de $\bb{F}_4$ es $\{1,a\}$ y $\eta$ deja fijo a $\bb{F}_2$, tenemos que:
    \begin{equation*}
        \eta(\alpha+\beta a) = \eta(\alpha)+\eta(\beta)\eta(a) = \alpha + \beta \eta(a)\qquad \forall \alpha,\beta\in \bb{F}_2
    \end{equation*}

    Por lo que $\eta$ está determinado por $\eta(a)$, y como $\eta(a)$ ha de ser raíz de $x^2+x+1$, tenemos que los homomorfismos de cuerpos $\eta:\bb{F}_4\to \bb{F}_{16}$ vienen dados por:
    \begin{align*}
        \eta_1(a) &= b^5 = b^2+b \\
        \eta_2(a) &= b^{10} = b^2+b+1
    \end{align*}




    \begin{comment}
    

    





    \noindent
    Cada $\eta$ por restricción nos da un homomorfismo de grupos $\eta:\bb{F}_4^\times\to \bb{F}_{16}^\times$ que, por el ejercicio anterior, tienen cardinales 3 y 15. Como 3 divide a 15, hay homomorfismos. Sabemos que $\bb{F}_2(a) = \langle a \rangle $ por ser 3 primo. Ahora, no estamos seguros de si $\bb{F}_2(b) = \langle b \rangle $, para lo cual hemos de probar que $O(b) = 15$.

    \begin{itemize}
        \item $b^2\neq 1$, ya que $b^2+1=0$, ya que $\{1,b,b^2, b^3\}$ es una $\bb{F}_2-$base de $\bb{F}_{16}$.
        \item $b^3\neq 1$ por la misma razón.
        \item $b^4 = b+1\neq 1$, ya que si no $b = 0$.
        \item $b^5 = b(b^4) = b(b+1) = b^2 + b \neq 1$, por la misma razón.
    \end{itemize}
    En definitiva, $O(b)=15$, luego $\bb{F}_{16}^{\times} = \langle b \rangle $.

    \noindent
    Buscando ahora homomorfismos de grupos, tenemos que llevar $a$ en un elemento de orden 3. Ahora, los candidatos a elementos de orden 3 de $\bb{F}_{16}^\times$ son los que generan un grupo de orden $5$ y $10$, es decir, $b^5$ y $b^{10}$, y tenemos que comprobar que son raíces de $x^2+x+1$.

    Finalmente, evalúo $p$ en las candidatas para comprobar que sean raíces:
    \begin{align*}
        p(b^5) &= b^{10} + b^{5} + 1 = {(b^2 + b)}^{2} + b^2 + b + 1 = b^4 + b^2 + b^2 + b + 1 \\
               &= b^4 + b + 1 = 0 
    \end{align*}
    Por el Automorfismo de Frobenius, la otra raíz es $b^{10}$. Sabemos que hay un $\eta$ para cada raíz del polinomio, obteniendo $\eta_i:\bb{F}_4\to \bb{F}_{16}$:
    \begin{equation*}
        \eta_1(a) = b^5, \qquad \eta_2(a) = b^{10}
    \end{equation*}
    \end{comment}
\end{ejercicio}


\begin{ejercicio}
    Demostrar que los anillos $\frac{\mathbb{Z}[i]}{\langle 3 \rangle }$ y $\frac{\bb{F}_3[x]}{\langle x^2+x+2 \rangle }$ son isomorfos sin necesidad de dar un isomorfismo concreto. ¿Serías capaz de darlo? ¿Y de calcularlos todos?\\

    \noindent
    Tenemos que $\frac{\bb{F}_3[x]}{\langle x^2+x+2 \rangle }$ es un cuerpo, porque $\bb{F}_3[x]$ es un DE y $x^2+x+2$ es irreducible por no tener raíces en $\bb{F}_3$. Además, al dividir entre $x^2+x+2$ obtenemos que los elementos de $\frac{\bb{F}_3[x]}{\langle x^2+x+2 \rangle }$ son las clases de equivalencia de los polinomios de grado menor que 2, es decir, los polinomios:
    \begin{equation*}
        \{a+bx : a,b\in \bb{F}_3\}\qquad x^2+x+2 = 0
    \end{equation*}

    Por tanto, $\frac{\bb{F}_3[x]}{\langle x^2+x+2 \rangle }$ tiene $3^2 = 9$ elementos. Por otra parte, $3$ es irreducible en el DE $\mathbb{Z}[i]$, ya que si $3 = zw$, entonces empleando la norma $N(z)$ al cuadrado, tenemos que:
    \begin{equation*}
        9 = |z|^2 |w|^2 
    \end{equation*}

    Como $|z|^2, |w|^2 \in \mathbb{Z}$ y son de la forma $a^2+b^2$ con $a,b\in \mathbb{Z}$, tenemos que $\{|z|^2, |w|^2\}= \{1,9\}$, luego o $z$ o $w$ son unidades, de donde $3$ es irreducible, de donde $\frac{\mathbb{Z}[i]}{\langle 3 \rangle }$ es un cuerpo. Veamos cuántos elementos tiene:
    \begin{description}
        \item[Opción 1:] Como $\langle 3 \rangle = \{3z : z\in \mathbb{Z}[i]\}$, tenemos que:
        $$z_1 + \langle 3 \rangle = z_2 + \langle 3 \rangle\iff z_1 - z_2 \in \langle 3 \rangle \iff z_1 - z_2 = 3z \text{ para algún } z\in \mathbb{Z}[i]$$

        Por tanto, dado $z_1 = a_1 + b_1 i$ y $z_2 = a_2 + b_2 i$, tenemos que:
        \begin{align*}
            z_1 + \langle 3 \rangle = z_2 + \langle 3 \rangle &\iff 3|\left(a_1 - a_2\right) \text{ y } 3|\left(b_1 - b_2\right)
        \end{align*}

        Por tanto, los elementos de $\frac{\mathbb{Z}[i]}{\langle 3 \rangle }$ son las clases de equivalencia de los elementos de $\mathbb{Z}[i]$ módulo $3$, es decir, los elementos:
        \begin{equation*}
            \{a+bi : a,b\in \{0,1,2\}\}\qquad i^2 = -1
        \end{equation*}
        Por lo que $\frac{\mathbb{Z}[i]}{\langle 3 \rangle }$ tiene $3^2 = 9$ elementos.

        \item[Opción 2:] Tenemos que:
        \begin{align*}
            \dfrac{\mathbb{Z}[i]}{\langle 3 \rangle } &= \left\{a+bi + \langle 3 \rangle : a,b\in \mathbb{Z}\right\}
            =\\&= \left\{(a+\langle 3 \rangle) + (b+\langle 3 \rangle)(i+\langle 3 \rangle) : a,b\in \mathbb{Z}\right\}
            =\\&= \left\{(a+\langle 3 \rangle) + (b+\langle 3 \rangle)(i+\langle 3 \rangle) : a,b\in \{0,1,2\}\right\}
        \end{align*}

        Por tanto, como máximo hay $3^2 = 9$ elementos; aunque podría tenermenos. Calculamos su subcuerpo primo haciendo uso de que la unidad de $\frac{\mathbb{Z}[i]}{\langle 3 \rangle }$ es $1+\langle 3 \rangle$ y sumándola consigo misma:
        \begin{align*}
            \left\{0+\langle 3 \rangle, 1+\langle 3 \rangle, 2+\langle 3 \rangle\right\}
        \end{align*}
        donde se ha empleado que $3+\langle 3 \rangle = 0+\langle 3 \rangle$. Por tanto, sabemos que $\frac{\mathbb{Z}[i]}{\langle 3 \rangle }$ tiene $3$ o $9$ elementos. Veamos que $i+\langle 3 \rangle$ no pertenece al subcuerpo primo. Si perteneciera, existiría $k\in \{0,1,2\}$ tal que:
        \begin{equation*}
            i+\langle 3 \rangle = k + \langle 3 \rangle \Longrightarrow i-k \in \langle 3 \rangle
        \end{equation*}

        Por tanto, $i-k = 3z$ para algún $z\in \mathbb{Z}[i]$, luego existen $a,b\in \mathbb{Z}$ tales que:
        \begin{equation*}
            i-k = 3(a+bi) \Longrightarrow 3a = -k \text{ y } 3b = 1
        \end{equation*}

        Como $3\notin \cc{U}(\mathbb{Z})$, no existe $b\in \mathbb{Z}$ tal que $3b = 1$, por lo que $i+\langle 3 \rangle$ no pertenece al subcuerpo primo.Por tanto, tenemos que $\frac{\mathbb{Z}[i]}{\langle 3 \rangle }$ tiene $9$ elementos.
    \end{description}

    \noindent
    Por tanto, $\frac{\mathbb{Z}[i]}{\langle 3 \rangle }$ y $\frac{\bb{F}_3[x]}{\langle x^2+x+2 \rangle }$ son cuerpos de $9$ elementos, por lo que son isomorfos.\\

    Buscamos ahora dar los homorfismos de cuerpos $\eta:\frac{\mathbb{Z}[i]}{\langle 3 \rangle }\to \frac{\bb{F}_3[x]}{\langle x^2+x+2 \rangle }$. Como la característica de ambos cuerpos es 3, tenemos que el cuerpo primo de ambos es $\bb{F}_3$, y por tanto, $\eta$ ha de dejar fijo a $\bb{F}_3$. Tenemos la siguiente situación:
    \begin{figure}[h]
        \centering
        \shorthandoff{""}
        \begin{tikzcd}
            \mathbb{F}_3 \arrow[r, "\iota", hook] \arrow[rd, "\iota", hook] & \frac{\bb{F}_3[x]}{\langle x^2+x+2 \rangle }                    \\
                                                                            & \frac{\mathbb{Z}[i]}{\langle 3 \rangle } \arrow[u, "\eta"']
        \end{tikzcd}
        \shorthandon{""}
    \end{figure}

    Como $\frac{\mathbb{Z}[i]}{\langle 3 \rangle } = \bb{F}_3(i)$, por la Proposición de Extensión tenemos que hay tantos homomorfismos de cuerpos $\eta:\frac{\mathbb{Z}[i]}{\langle 3 \rangle }\to \frac{\bb{F}_3[x]}{\langle x^2+x+2 \rangle }$ como raíces tenga el polinomio dado por $x^2+1\in \frac{\bb{F}_3[x]}{\langle x^2+x+2 \rangle }[x]$.
    \begin{itemize}
        \item Sea $y=u+vx \in \frac{\bb{F}_3[x]}{\langle x^2+x+2 \rangle }$, con $u,v\in \bb{F}_3$. Entonces:
            \begin{align*}
                y^2 + 1 &= {(u+vx )}^2 + 1 = u^2 + 2uvx + v^2x^2 + 1 \\ &= u^2 + 1 + 2uvx + v^2(-x-2) = u^2 +1+2uvx + v^2(2x+1)
                =\\&= (u^2+1+v^2) + 2(uv+v^2)x
            \end{align*}
            Por lo que $y$ es raíz de $x^2+1$ si, y solo si:
            \begin{equation*}
                u^2+1+v^2 = 0 \qquad\text{y}\qquad 2v(u+v) = 0
            \end{equation*}
            de donde $v=0$ o $u+v=0$.
            \begin{itemize}
                \item Si $v=0$, tenemos que $u^2+1 = 0$, de donde $u^2 = 2$, y como $2$ no es un cuadrado en $\bb{F}_3$, no hay solución.
                \item Si $u+v=0$, como $u^2+1+v^2 = 0$, tenemos:
                \begin{align*}
                    u^2+1+v^2 = 0 &\Longleftrightarrow u^2+1+u^2 = 0 \Longleftrightarrow 2u^2 + 1 = 0 \Longleftrightarrow 2u^2 = 2 \Longleftrightarrow u^2 = 1 \\ &\Longrightarrow u = \pm 1 \Longrightarrow v = -u = \mp 1
                \end{align*}

                Por tanto, las raíces de $x^2+1$ en $\frac{\bb{F}_3[x]}{\langle x^2+x+2 \rangle }$ son:
                \begin{equation*}
                    1+2x \qquad\text{y}\qquad 2+x
                \end{equation*}
            \end{itemize}        
    \end{itemize}



    Por tanto, los homomorfismos de cuerpos $\eta:\frac{\mathbb{Z}[i]}{\langle 3 \rangle }\to \frac{\bb{F}_3[x]}{\langle x^2+x+2 \rangle }$ vienen dados por:
    \begin{align*}
        \eta_1(i) &= 1+2x \\
        \eta_2(i) &= 2+x
    \end{align*}

    Además, dar la imagen de $i$ determina completamente a $\eta$, ya que $\{1,i\}$ es una base de $\frac{\mathbb{Z}[i]}{\langle 3 \rangle }$ sobre $\bb{F}_3$, luego:
    \begin{align*}
        \eta(\alpha+\beta i) &= \eta(\alpha)+\eta(\beta)\eta(i) = \alpha + \beta \eta(i)\qquad \forall \alpha,\beta\in \bb{F}_3
    \end{align*}

    Por tanto, estos son los homomorfismos de cuerpos entre $\frac{\mathbb{Z}[i]}{\langle 3 \rangle }$ y $\frac{\bb{F}_3[x]}{\langle x^2+x+2 \rangle }$. Como todo homomorfismo es inyectivo, y ambos cuerpos tienen el mismo cardinal, tenemos que estos homomorfismos son isomorfismos de cuerpos.
\end{ejercicio}

\begin{ejercicio}
    Se pide:
    \begin{enumerate}
        \item Comprobar que $\sqrt{3}\in \mathbb{Q}\left(\sqrt{1+2\sqrt{3}}\right)$.

            Llamamos $\alpha = \sqrt{1+2\sqrt{3}}$ y calculamos:
            \begin{equation*}
                \alpha^2 = 1+2\sqrt{3} \quad \Longrightarrow \quad  \sqrt{3} = \dfrac{\alpha^{2}-1}{2}\in \mathbb{Q}(\alpha)
            \end{equation*}
            De donde deducimos que $\mathbb{Q}(\sqrt{3})\leq \mathbb{Q}(\alpha)$.
        \item Calcular $\Irr\left(\alpha, \mathbb{Q}\left(\sqrt{3}\right)\right)$.

            Sabemos que $\alpha$ es raíz de $f=x^2-1-2\sqrt{3}\in \mathbb{Q}\left(\sqrt{3}\right)[x]$, con lo que:
            \begin{equation*}
                \left[\mathbb{Q}(\alpha):\mathbb{Q}\left(\sqrt{3}\right)\right] \leq 2
            \end{equation*}
            Supongamos que $\left[\mathbb{Q}(\alpha):\mathbb{Q}\left(\sqrt{3}\right)\right]=1$, con lo que $\alpha\in \mathbb{Q}\left(\sqrt{3}\right)$. Como se tiene que $\Irr(\sqrt{3},\mathbb{Q}) = x^2-3$, tenemos que una base de $\mathbb{Q}\left(\sqrt{3}\right)$ sobre $\mathbb{Q}$ es $\{1,\sqrt{3}\}$, de donde $\alpha = a + b\sqrt{3}$ para ciertos $a,b\in \mathbb{Q}$. Si elevamos al cuadrado:
            \begin{equation*}
                1 + 2\sqrt{3} = \alpha^2 = a^2 + 3b^2 + 2ab\sqrt{3}
            \end{equation*}
            Usando que $\{1,\sqrt{3}\}$ es una base de $\mathbb{Q}\left(\sqrt{3}\right)$, tenemos entonces que:
            \begin{align*}
                \left.\begin{array}{l}
                    1 = a^2 + 3b^2 \\
                    2 = 2 ab
                    \end{array}\right\} &\Longrightarrow \left\{\begin{array}{l}
                    b = \nicefrac{1}{a} \\
                    1 = a^2 + 3\frac{1}{a^2}
            \end{array}\right\} \Longrightarrow a^2 = a^4 + 3\Longrightarrow a^4 - a^2 + 3 = 0
            \end{align*}

            Los únicos posibles racionales son $\pm 1$ y $\pm 3$, que no son raíces del polinomio, por lo que no hay solución en $\mathbb{Q}$, luego $\alpha\notin \mathbb{Q}\left(\sqrt{3}\right)$. Por tanto, $\left[\mathbb{Q}(\alpha):\mathbb{Q}\left(\sqrt{3}\right)\right]\neq 1$ y:
            \begin{equation*}
                \left[\mathbb{Q}(\alpha):\mathbb{Q}\left(\sqrt{3}\right)\right] = 2
            \end{equation*}
            Por tanto, $\Irr\left(\alpha,\mathbb{Q}\left(\sqrt{3}\right)\right) = x^2-1-2\sqrt{3}$.
            
        \item Calcular los homomofismos de $\mathbb{Q}(\alpha)$ en $\mathbb{C}$.
        
            Como la característica de ambos cuerpos es 0, tenemos que $\mathbb{Q}$ es el cuerpo primo de ambos, y por tanto, cualquier homomorfismo de cuerpos $\eta:\mathbb{Q}(\alpha)\to \mathbb{C}$ ha de dejar fijo a $\mathbb{Q}$. Por tanto, tenemos la siguiente situación:
            \begin{figure}[H]
                \centering
                \shorthandoff{""}
                \begin{tikzcd}
                \mathbb{Q} \arrow[rd, "\iota"', hook] \arrow[r, "\iota", hook] & \mathbb{C}                            \\
                                                                  & \mathbb{Q}(\alpha) \arrow[u, "\eta"']
                \end{tikzcd}
                \shorthandon{""}
            \end{figure}
            donde $\iota$ es la inclusión. Queremos por tanto calcular $\Ex(\iota,\iota)$.\\

            En vistas de aplicar la Proposición de extensión, necesitamos conocer un polinomio irreducible sobre $\mathbb{Q}$ que tenga a $\alpha$ como raíz. Puesto que no es directo, empleamos que:
            $$\bb{Q}\leq \bb{Q}(\sqrt{3})\leq \bb{Q}(\alpha)$$
            Por tanto, tenemos la siguiente situación:
            \begin{figure}[h]
                \centering
                \shorthandoff{""}
                \begin{tikzcd}
                    \mathbb{Q} \arrow[r, "\iota", hook] \arrow[rd, "\iota", hook] \arrow[d, "\iota", hook] & \mathbb{C}                            \\
                    \mathbb{Q}(\sqrt{3}) \arrow[r, "\iota", hook] & \mathbb{Q}(\alpha) \arrow[u, "\eta"']
                \end{tikzcd}
                \shorthandon{""}
            \end{figure}

            Por tanto, tenemos que:
            \begin{equation*}
                \Ex(\iota,\iota\iota) = \biguplus_{\theta\in \Ex(\iota,\iota)} \Ex(\theta,\iota)
            \end{equation*}


            Calculamos primero $\Ex(\iota,\iota)$, que es el conjunto de homomorfismos de cuerpos $\mathbb{Q}(\sqrt{3})\to \mathbb{C}$ que dejan fijo a $\mathbb{Q}$. Sabemos que $\Irr(\sqrt{3},\mathbb{Q}) = x^2-3$, con raíces $\pm \sqrt{3}$, por lo por la Proposición de extensión, tenemos que:
            \begin{equation*}
                \Ex(\iota,\iota) = \{\eta_0, \eta_1\}
            \end{equation*}
            donde $\eta_j$ está determinada por:
            \begin{equation*}
                \eta_j(\sqrt{3}) = {(-1)}^{j}\sqrt{3}, \qquad \forall j\in \{0,1\}
            \end{equation*}
            Por tanto, tenemos que:
            \begin{equation*}
                \Ex(\iota,\iota\iota) = \Ex(\eta_0,\iota) \biguplus \Ex(\eta_1,\iota)
            \end{equation*}
            \begin{itemize}
                \item Para $\theta = \eta_0$, tenemos que $\Irr(\alpha,\mathbb{Q}(\sqrt{3})) = x^2-1-2\sqrt{3}$, con raíces $\pm \alpha$, por lo que:
                    \begin{equation*}
                        \Ex(\eta_0,\iota) = \{\eta_{0,0},\eta_{0,1}\}
                    \end{equation*}
                    donde $\eta_{0,k}$ está determinada por:
                    \begin{equation*}
                        \eta_{0,k}(\alpha) = {(-1)}^{k}\alpha, \qquad \forall k\in \{0,1\}
                    \end{equation*}
                \item Para $\theta = \eta_1$, tenemos que $\Irr(\alpha,\mathbb{Q}(\sqrt{3}))^{\eta_1} = x^2-1+2\sqrt{3}$, con raíces $\pm \beta$, donde $\beta = \sqrt{1-2\sqrt{3}}$, por lo que:
                    \begin{equation*}
                        \Ex(\eta_1,\iota) = \{\eta_{1,0},\eta_{1,1}\}
                    \end{equation*}
                    donde $\eta_{1,k}$ está determinada por:
                    \begin{equation*}
                        \eta_{1,k}(\alpha) = {(-1)}^{k}\beta, \qquad \forall k\in \{0,1\}
                    \end{equation*}
            \end{itemize}
            
            Por tanto, tenemos que:
            \begin{equation*}
                \Ex(\iota,\iota\iota) = \{\eta_{0,0},\eta_{0,1},\eta_{1,0},\eta_{1,1}\}
            \end{equation*}


        \item Calcular $\Irr(\alpha,\mathbb{Q})$ y sus raíces en $\mathbb{C}$.
            
            Empleando el Lema de la torre, tenemos que:
            \begin{equation*}
                [\mathbb{Q}(\alpha):\mathbb{Q}] = [\mathbb{Q}(\alpha):\mathbb{Q}\left(\sqrt{3}\right)]\cdot [\mathbb{Q}\left(\sqrt{3}\right):\mathbb{Q}] = 2\cdot 2 = 4
            \end{equation*}
            Por tanto, $\Irr(\alpha,\mathbb{Q})$ es un polinomio de grado 4. Para calcularlo, partimos de:
            \begin{align*}
                \alpha^2-1 &= 2\sqrt{3} \Longrightarrow \alpha^4 - 2\alpha^2 + 1 = 12 \Longrightarrow \alpha^4 - 2\alpha^2 - 11 = 0
            \end{align*}

            Por tanto, $\Irr(\alpha,\mathbb{Q}) = x^4-2x^2-11$. Resolvemos como una cuadrática en $x^2$, que sabemos hacerlo pues $\car(\bb{Q})\neq 2$:
            \begin{align*}
                x^2 &= \frac{2\pm \sqrt{4+44}}{2} = 1\pm \sqrt{12} = 1\pm 2\sqrt{3}
            \end{align*}
            Por tanto, las raíces de $\Irr(\alpha,\mathbb{Q})$ en $\mathbb{C}$ son:
            \begin{equation*}
                \pm \sqrt{1+2\sqrt{3}} = \pm \alpha \qquad\text{y}\qquad \pm \sqrt{1-2\sqrt{3}} = \pm \beta
            \end{equation*}
    \end{enumerate}
\end{ejercicio}


\begin{ejercicio}
    ¿Son los cuerpos $\mathbb{Q}(\sqrt{2})$ y $\mathbb{Q}(\sqrt{3})$ isomorfos?\\

    Puesto que el cuerpo $\mathbb{Q}$ es el cuerpo primo de ambos, cualquier homomorfismo de cuerpos $\eta:\mathbb{Q}(\sqrt{2})\to \mathbb{Q}(\sqrt{3})$ ha de dejar fijo a $\mathbb{Q}$. Por tanto, tenemos la siguiente situación:
    \begin{figure}[H]
        \centering
        \shorthandoff{""}
        \begin{tikzcd}
            \mathbb{Q} \arrow[rd, "\iota"', hook] \arrow[r, "\iota", hook] & \mathbb{Q}(\sqrt{3})                            \\
                                                              & \mathbb{Q}(\sqrt{2}) \arrow[u, "\eta"']
        \end{tikzcd}
        \shorthandon{""}
    \end{figure}

    Consideramos $\Irr(\sqrt{2},\mathbb{Q}) = x^2-2$. Por la Proposición de extensión, tenemos que $\Ex(\iota,\iota)$ es biyectivo con el conjunto de raíces de $x^2-2$ en $\mathbb{Q}(\sqrt{3})$. Por tanto, para ver si existen homomorfismos de cuerpos $\eta:\mathbb{Q}(\sqrt{2})\to \mathbb{Q}(\sqrt{3})$, basta con ver si $x^2-2$ tiene raíces en $\mathbb{Q}(\sqrt{3})$.\\

    Como $\{1,\sqrt{3}\}$ es una base de $\mathbb{Q}(\sqrt{3})$ sobre $\mathbb{Q}$, tenemos que $x^2-2$ tiene raíces en $\mathbb{Q}(\sqrt{3})$ si, y solo si, existen $a,b\in \mathbb{Q}$ tales que:
    \begin{equation*}
        (a+b\sqrt{3})^2 - 2 = 0 \Longrightarrow a^2 + 3b^2 + 2ab\sqrt{3} - 2 = 0 \Longrightarrow a^2 + 3b^2 - 2 = 0 \quad\text{y}\quad 2ab = 0
    \end{equation*}
    Por tanto, $a=0$ o $b=0$.
    \begin{itemize}
        \item Si $a=0$, tenemos que $3b^2-2 = 0$. Como $b$ es una raíz racional de ese polinomio, tenemos que $b\in \{\pm 1, \pm 2, \pm \nicefrac{1}{3}, \pm \nicefrac{2}{3}\}$, y ninguno de esos valores es raíz del polinomio, por lo que no hay solución en $\mathbb{Q}$.
        \item Si $b=0$, tenemos que $a^2-2 = 0$, de donde $a^2 = 2$, que no tiene solución en $\mathbb{Q}$.
    \end{itemize}

    Por tanto, $x^2-2$ no tiene raíces en $\mathbb{Q}(\sqrt{3})$, y por tanto, no existen homomorfismos de cuerpos $\eta:\mathbb{Q}(\sqrt{2})\to \mathbb{Q}(\sqrt{3})$. En particular, no existen isomorfismos de cuerpos entre $\mathbb{Q}(\sqrt{2})$ y $\mathbb{Q}(\sqrt{3})$.
\end{ejercicio}